\documentclass[11pt]{ctexart}
\usepackage[a4paper,margin=1in]{geometry}
\usepackage{graphicx}
\usepackage{xcolor}
\usepackage{hyperref}
\definecolor{refgray}{gray}{0.45}
\title{\textbf{Market Views｜全球投行叙事汇编}\\[0.4em]{\large 宏观情绪极值分化，AI结构性短缺与盈利兑现成为市场主线}}
\author{KC桌面 自动生成}
\date{260715}
\begin{document}
\maketitle
\tableofcontents
\newpage
\section{一页摘要}
\begin{itemize}
  \item 全球宏观情绪从谨慎转向高度乐观，市场预期经济‘不着陆’比例创54\%历史新高，硬着陆预期降至2\%，但驱动力正从估值重估转向盈利兑现。
  \item AI相关交易成为最拥挤方向，做多全球半导体被82\%基金经理视为最拥挤交易，资金集中度与估值分化达到极端水平，脆弱性上升。
  \item 美元多头仓位仍在高位但增速放缓，日元空头出现首个去杠杆信号，人民币中间价模型显示逆周期因子使用力度减弱，定价机制向市场化倾斜。
  \item AI基础设施需求持续超预期，先进制程、存储及封装产能全面紧张，结构性短缺取代周期性波动成为行业主线，台积电产能瓶颈成为公开约束条件。
  \item AI基础设施投资正从概念验证转向规模部署，但受益结构出现分化：算力层需求强劲，平台层和内容层的差异化价值开始被重新定价。
  \item 全球汽车与零部件行业结构性分化：欧洲传统车企受中国需求下滑与成本压力挤压利润率，中国新能源品牌凭借产品换代与品牌溢价实现逆势增长。
  \item 全球消费市场显著分化：西欧美妆个护销量驱动增长，奢侈品板块预期已充分定价疲软，中国宠物食品与零食行业出现结构性改善信号。
  \item 支付处理商低PE估值存在系统性误读，非GAAP盈利因持续计入重组与转型费用而高估真实盈利能力，基于‘清洁盈利’框架的真实估值显著高于表面倍数。
\end{itemize}
\section{全球投行叙事汇编}
今日纳入 68 篇投行/券商研究，按 12 个市场、资产与产业主题系统整理。
\newpage
\subsection{宏观与策略：情绪极值与结构分化}
\textbf{全球宏观情绪已从谨慎转向高度乐观，但市场驱动力正从估值重估转向盈利兑现，资金集中度与估值分化达到极端水平，脆弱性上升。}
\subsubsection*{主流共识}
\begin{itemize}
  \item 全球资金持续流入美国股票基金，7月节奏与6月基本持平，美国资产仍是相对主线。
  \item 市场预期经济‘不着陆’的比例创历史新高（54\%），硬着陆预期降至2\%的历史低点。
  \item AI相关交易成为最拥挤方向，做多全球半导体被82\%的基金经理视为最拥挤交易。
\end{itemize}
\subsubsection*{分歧与条件差异}
\begin{itemize}
  \item 通胀性质分歧：MS认为当前通胀是成本传导与供给扰动，并非需求回暖驱动的再通胀；而市场定价隐含对再通胀的担忧。
  \item 中国资产风格分歧：JPM指出昂贵领涨结构（AI硬件）估值比达18倍（114.9x vs 6.3x），7月初已出现动量/成长向价值/低波动的风格反转；GS则观察到新房成交环比走弱但同比仍正，政策重心转向长期制度框架。
  \item 亚太科技方向分歧：MS强调AI对产业的影响正从概念走向业绩兑现（联想EPS比共识高20\%），但JPM认为AI交易已拥挤且脆弱。
\end{itemize}
\subsubsection*{机构观点}
\paragraph{BofA} 7月全球基金经理调查显示现金仓位从4.1\%降至3.6\%，触发BofA反向观察信号；预期‘不着陆’的投资者比例创54\%历史新高，硬着陆预期仅2\%。牛熊指标升至9.4，进入极端看多区间。历史规律显示现金≤3.6\%后全球市场短期承压，美债相对稳健。仓位集中、预期一致，反向指标亮红灯。
{\small\color{refgray}数据：现金仓位从4.1\%降至3.6\%，触发BofA反向观察信号\par}
{\small\color{refgray}数据：预期‘不着陆’比例54\%，创历史新高\par}
{\small\color{refgray}数据：硬着陆预期降至2\%，历史低点\par}
{\small\color{refgray}数据：牛熊指标升至9.4，进入极端看多区间\par}
{\small\color{refgray}边际变化：现金仓位降至3.6\%触发反向观察信号，此前数月处于4\%以上；牛熊指标从8.0附近升至9.4，情绪从乐观升级为极端看多。\par}
\paragraph{MS} 当前通胀是‘通胀而不再通胀’——价格上升来自成本传导与供给扰动，而非需求回暖。下游利润被压缩，消费者信心偏弱，储蓄倾向仍高。财政重点在加快预算执行而非扩大规模，下半年约2万亿元财政脉冲空间，但出口韧性降低了加码紧迫性。央行按兵不动，政策基调未变。
{\small\color{refgray}数据：剔除黄金后核心CPI同比涨幅明显低于整体CPI\par}
{\small\color{refgray}数据：下半年约2万亿元财政脉冲空间\par}
{\small\color{refgray}数据：出口韧性降低了追加逆周期调节的紧迫性\par}
{\small\color{refgray}数据：消费者信心偏弱，储蓄倾向仍高\par}
{\small\color{refgray}边际变化：市场此前对下半年财政扩张有较高期待，但MS指出重点在加快执行而非扩大规模，政策节奏比预期更平稳。\par}
\paragraph{NOM} 7月前两周跨境资金显示全球资金持续流入美国股票基金，7月前9天流入44亿美元，与6月同期49亿美元节奏基本持平。债券基金流入仅5900万美元，远低于股票。新兴市场ETF流入回暖但力度有限。韩国个人投资者从净买入转为净卖出美国资产8400万美元，台湾投资者继续净卖出美元债券基金但速度放缓。
{\small\color{refgray}数据：7月前9天美国股票基金净流入44亿美元\par}
{\small\color{refgray}数据：6月全月美国股票基金净流入136亿美元\par}
{\small\color{refgray}数据：7月第一周债券基金净流入仅5900万美元\par}
{\small\color{refgray}数据：韩国个人投资者净卖出美国资产8400万美元（此前一周净买入2100万）\par}
{\small\color{refgray}边际变化：韩国个人投资者从连续三周净买入转为净卖出，台湾投资者减持美元债券基金速度放缓，亚洲零售资金行为出现分化。\par}
\paragraph{GS} 全球市场分化明显，美国与能源板块录得正收益，欧洲和新兴市场承压。市场驱动力正从估值重估转向盈利兑现，HALO交易（重资产、低过时率）进入第二阶段。能源板块领涨不仅是地缘事件驱动，更反映结构性资金再配置——高实际利率、地缘碎片化和供应链重构强化了对有形生产性资产的偏好。
{\small\color{refgray}数据：MSCI全球指数上周基本持平\par}
{\small\color{refgray}数据：欧洲斯托克600下跌1.8\%，德国DAX下跌2.8\%\par}
{\small\color{refgray}数据：布伦特原油价格随中东局势升温走高\par}
{\small\color{refgray}数据：GS不确定性偏好指标（RAI）从6月高点回落但仍处历史较高分位\par}
{\small\color{refgray}边际变化：市场驱动力从估值重估转向盈利兑现，赢家与输家分化将加剧；能源板块上涨不仅是地缘事件驱动，更是结构性资金再配置的缩影。\par}
\paragraph{JPM} 6月MSCI中国指数下跌7.1\%，但动量因子上涨9.8\%、成长因子上涨7.4\%，涨幅高度集中在少数AI硬件公司。昂贵组市盈率中位数114.9倍，便宜组仅6.3倍，估值比超18倍，与2020年结构不同（当前集中于科技硬件、半导体、材料）。7月初风格反转：动量跌4.7\%、成长跌4.5\%，价值反弹7.4\%、低波动涨8\%。
{\small\color{refgray}数据：6月MSCI中国指数下跌7.1\%，73\%公司收跌\par}
{\small\color{refgray}数据：动量因子上涨9.8\%，成长因子上涨7.4\%\par}
{\small\color{refgray}数据：昂贵组市盈率中位数114.9倍，便宜组6.3倍\par}
{\small\color{refgray}数据：7月初动量因子下跌4.7\%，成长因子下跌4.5\%\par}
{\small\color{refgray}边际变化：7月初市场风格从动量/成长急剧转向价值/低波动，此前6月动量与成长大幅跑赢；估值分化结构从2020年生物医药转向AI硬件，集中度更高。\par}
\paragraph{MS} 亚太地区三个可观察方向：1）Keyence受益于半导体设备短期需求与物理AI长期主题叠加；2）联想AI驱动定价能力增强，2027-2029财年EPS比市场共识高约20\%；3）智邦科技2Q营收环比增速从预期20\%上调至25-30\%，下半年动能有望持续。AI对产业的影响正从概念走向可量化的业绩兑现。
{\small\color{refgray}数据：联想2027-2029财年EPS比市场共识高约20\%\par}
{\small\color{refgray}数据：智邦科技2Q营收环比增速从20\%上调至25-30\%\par}
{\small\color{refgray}数据：Keyence同时受益于半导体设备短期需求与物理AI长期主题\par}
{\small\color{refgray}边际变化：AI对产业的影响从概念走向可量化的业绩兑现，联想定价能力增强、智邦营收超预期，市场对自动化设备的关注从周期性订单转向结构性需求。\par}
\paragraph{GS} 7月第三周约75城新房成交面积环比下降25\%，同比转负6\%，但7月累计新房同比仍正5\%，二手房同比正11\%。二手房带看和认购环比分别下降8\%和6\%，价格预期出现边际下调。库存月数降至27.0个月，低于6月均值27.5个月。‘十五五’扩大消费规划首次将房地产纳入‘增加居民财产性收入’框架，但未提及专项债收储和税费减免等短期工具。
{\small\color{refgray}数据：7月第三周新房成交环比下降25\%，同比-6\%\par}
{\small\color{refgray}数据：7月累计新房同比+5\%，二手房同比+11\%\par}
{\small\color{refgray}数据：二手房带看和认购环比分别下降8\%和6\%\par}
{\small\color{refgray}数据：库存月数27.0个月，低于6月均值27.5个月\par}
{\small\color{refgray}边际变化：7月成交环比走弱但同比仍正，政策重心从短期支持转向长期制度框架（房地产纳入财产性收入），短期工具缺席。\par}
\subsubsection*{关键数据}
\begin{itemize}
  \item 现金仓位降至3.6\%，触发BofA反向观察信号
  \item 预期‘不着陆’比例54\%，创历史新高
  \item 昂贵组市盈率中位数114.9倍 vs 便宜组6.3倍
  \item 7月前9天美国股票基金净流入44亿美元
  \item 联想2027-2029财年EPS比共识高约20\%
  \item 7月第三周新房成交环比下降25\%
  \item HALO交易第一阶段（估值重估）已基本完成
\end{itemize}
\begin{figure}[htbp]
\centering
\includegraphics[width=0.88\linewidth]{figures/fig_121.jpg}
\caption{Exhibit 1 - Exhibit 1: Primary GFA sold last week was \$-25\textbackslash{}\%\$ wow and \$-6\textbackslash{}\%\$ yoy in c.75 cities Exhibit 2: Primary GFA sold YTD on average was -12\% yoy in c.75 cities, and -13\%/-38\% vs. 2024/2023 level Secondary market: Week 28/YTD volumes}
\end{figure}
\begin{figure}[htbp]
\centering
\includegraphics[width=0.88\linewidth]{figures/fig_136.jpg}
\caption{Exhibit 1 - Exhibit 1: Global market performance MSCI Indices, 1-week price return (USD, \%) Exhibit 2: World equity indices USD, indexed price performance Exhibit 3: MSCI AC World sector performance MSCI Indices, 1-week price return (USD,}
\end{figure}
\begin{figure}[htbp]
\centering
\includegraphics[width=0.88\linewidth]{figures/fig_161.jpg}
\caption{Figure 1 - Figure 1: Valuation Ratio: Expensive vs. Cheap That does not automatically make the expensive stocks wrong, nor the cheap stocks right. Expensive can stay expensive when earnings momentum, liquidity, and narrative all point in t}
\end{figure}
\begin{figure}[htbp]
\centering
\includegraphics[width=0.88\linewidth]{figures/fig_183.jpg}
\caption{Exhibit 2 - Exhibit 2: Cumulative outperformance Exhibit 3: Weekly hit ratio Three Actionable Ideas are not and should not be considered a portfolio: Each investment idea is chosen based on its own merit and without any consideration of th}
\end{figure}
{\small\color{refgray}本节综合 7 篇研究；涉及 BofA、MS、NOM、GS、JPM。}
\newpage
\subsection{外汇仓位与人民币定价：美元多头趋稳，日元空头初现回补}
\textbf{美元多头仓位仍在高位但增速放缓，日元空头出现首个去杠杆信号，人民币中间价模型显示逆周期因子使用力度减弱，定价机制向市场化倾斜。}
\subsubsection*{主流共识}
\begin{itemize}
  \item 美元多头仓位仍占主导，但新增动能减弱，投机性期货仓位从51.2\%升至52.6\%，增速低于前几周。
  \item 日元空头回补初现，但整体仍偏空，仓位从+1.5σ降至+1.1σ，结构性贬值趋势尚未逆转。
  \item 人民币贬值预期降温，USD/CNY看跌期权需求回落，交易量低于2022-2024年水平。
\end{itemize}
\subsubsection*{分歧与条件差异}
\begin{itemize}
  \item 期权与期货市场对日元仓位的反应速度不一致：期权市场录得2.5σ单日净看跌需求，但期货空头回补在消息前已开始，显示调整可能接近尾声而非起点。
  \item 人民币中间价模型显示逆周期因子偏差仅23bp，远低于历史区间，暗示政策干预意愿减弱；但JPM仓位数据仍显示USD/CNY看跌需求未完全翻多，市场情绪分歧。
  \item G10货币中，挪威克朗和日元仓位变化与整体美元偏强格局形成对比，MS反向策略建议做多NOK/JPY，而JPM认为日元空头回补持续性存疑。
\end{itemize}
\subsubsection*{机构观点}
\paragraph{MS} 截至7月10日，期权市场显示战术投资者整体净做多美元、做空欧元，投机性美元期货仓位升至未平仓合约的52.6\%（前周51.2\%），增速放缓。挪威克朗兑欧元多头和瑞郎空头在期权市场增加，日元期货多头和加元空头在期货市场上升。资产管理人偏多欧元、空纽元，杠杆基金多澳元、空加元，机构内部分歧较大。基于仓位百分位数的反向策略建议做多NOK/JPY。
{\small\color{refgray}数据：美元期货仓位从51.2\%升至52.6\%\par}
{\small\color{refgray}数据：G10中纽元情绪改善最明显\par}
{\small\color{refgray}数据：期权市场显示投资者增加NOK多头、CHF空头\par}
{\small\color{refgray}数据：期货市场增加JPY多头、CAD空头\par}
{\small\color{refgray}边际变化：美元多头仓位增速放缓，新增多头力量减弱；挪威克朗和日元成为仓位变化中的异类，与整体美元偏强格局形成对比。\par}
\paragraph{NOM} 美元兑人民币中间价模型预测值从6.7972下调至6.7914，调整约58bp。加入逆周期因子后预测值为6.7949，与不含因子结果仅差23bp，偏差收窄至历史低位，表明逆周期因子使用力度减弱，定价机制向市场化倾斜。模型显示短期中间价在6.79-6.80区间窄幅波动，后续关注7月下旬重要会议和8月央行Q2货币政策报告。
{\small\color{refgray}数据：模型预测值6.7914，前次6.7972\par}
{\small\color{refgray}数据：含逆周期因子预测值6.7949，偏差23bp\par}
{\small\color{refgray}数据：相比上次官方收盘价上调97bp\par}
{\small\color{refgray}数据：相比上次中间价下调23bp\par}
{\small\color{refgray}边际变化：逆周期因子偏差从历史高位收窄至23bp，暗示政策层对市场定价的干预意愿减弱，认为当前汇率水平接近合意区间。\par}
\paragraph{JPM} 日元空头出现首个去杠杆信号：上周五USD/JPY期权市场录得自2025年8月以来最大单日净看跌需求，量级达2.5σ，触发因素为GPIF可能调整外汇对冲策略的言论。综合期权与期货数据，USD/JPY多头仓位从+1.5σ降至+1.1σ，但仍偏空。人民币贬值预期降温，USD/CNY看跌期权需求明显回落，交易量低于2022-2024年水平，但尚未完全翻多。韩元和台币仍被做空，与人民币形成对比。
{\small\color{refgray}数据：USD/JPY期权单日净看跌需求2.5σ\par}
{\small\color{refgray}数据：USD/JPY多头仓位从+1.5σ降至+1.1σ\par}
{\small\color{refgray}数据：期权行权价中位数159.75，10\%分位156.0\par}
{\small\color{refgray}数据：USD/CNY看跌期权需求回落，交易量低于2022-2024年\par}
{\small\color{refgray}边际变化：日元空头回补初现，但报告怀疑能否逆转结构性贬值趋势；期货市场空头回补在消息前已开始，显示调整可能接近尾声而非起点。人民币贬值预期降温，但韩元台币做空仍在增加。\par}
\subsubsection*{关键数据}
\begin{itemize}
  \item 美元期货仓位从51.2\%升至52.6\%，增速放缓
  \item USD/JPY期权单日净看跌需求2.5σ，为2025年8月以来最大
  \item USD/JPY多头仓位从+1.5σ降至+1.1σ
  \item 人民币中间价模型预测6.7914，逆周期因子偏差仅23bp
  \item USD/CNY看跌期权需求回落，交易量低于2022-2024年
  \item 期权行权价中位数159.75，10\%分位156.0
  \item G10中纽元情绪改善最明显
  \item 挪威克朗多头和瑞郎空头在期权市场增加
\end{itemize}
\begin{figure}[htbp]
\centering
\includegraphics[width=0.88\linewidth]{figures/fig_166.jpg}
\caption{Figure 3 - Figure 3: USD/JPY longs are not stretched and are off the highs - but are still material}
\end{figure}
\begin{figure}[htbp]
\centering
\includegraphics[width=0.88\linewidth]{figures/fig_169.jpg}
\caption{Figure 3 - Figure 5: USD/CNY downside demand has cooled... USD/CNY net options flow (calls less puts, notionals) (z-score)}
\end{figure}
\begin{figure}[htbp]
\centering
\includegraphics[width=0.88\linewidth]{figures/fig_040.jpg}
\caption{Exhibit 1 - Exhibit 1: Molly Nickolin Strategist}
\end{figure}
\begin{figure}[htbp]
\centering
\includegraphics[width=0.88\linewidth]{figures/fig_041.jpg}
\caption{Exhibit 2 - Exhibit 2: In the week ending Friday, July 10, options pricing data indicate that investors increased long NOK (vs. EUR) positions and increased short CHF positions (Exhibit 3). In the futures market, investors increased long JPY p}
\end{figure}
{\small\color{refgray}本节综合 3 篇研究；涉及 MS、NOM、JPM。}
\newpage
\subsection{半导体与设备：AI驱动结构性短缺，产能瓶颈重塑竞争格局}
\textbf{AI基础设施需求持续超预期，导致先进制程、存储及封装产能全面紧张，结构性短缺正取代周期性波动成为行业主线。台积电产能瓶颈迫使客户寻求替代方案，但短期难以撼动其领先地位；利基型DRAM上行周期可能延续至2028年；模拟芯片进入可持续盈利修复早期阶段；PCB与封装设备环节受益于AI扩产，竞争焦点从进入供应链转向规模化交付能力。}
\subsubsection*{主流共识}
\begin{itemize}
  \item AI基础设施需求持续拉动半导体上游，先进制程、HBM、先进封装产能全面紧张，结构性短缺取代周期性波动。
  \item 台积电产能瓶颈成为行业公开约束条件，客户开始寻找三星、英特尔等替代方案，但台积电领先至少一个世代，短期影响有限。
  \item DRAM供需缺口预计持续至2028年，DDR5/DDR4价格下半年继续走高，服务器收入占比提升改善产品结构。
\end{itemize}
\subsubsection*{分歧与条件差异}
\begin{itemize}
  \item 市场预期台积电3Q26毛利率达69\%-70\%，但MS认为海外厂和2nm稀释效应将压制在67\%-68\%，分歧约2个百分点。
  \item 模拟芯片订单改善是库存回补还是终端需求真实加速？MS基于定价纪律、订单取消量有限和客户可见度拉长，判断为后者。
  \item 联发科获得谷歌TPU部分订单，市场担忧博通份额被结构性取代，但MS认为从参与到规模化替换存在巨大执行鸿沟，博通仍将维持约80\%份额。
  \item 内存TAM预期已跑在CSP实际资本开支前面，短期情绪取决于CSP二季度业绩披露后资本开支指引上调幅度。
\end{itemize}
\subsubsection*{机构观点}
\paragraph{Bernstein} 台积电2Q26营收NT\$1.27T，环比+12\%，同比+36\%，落在指引高端，比市场共识高约0.5\%。产能瓶颈已导致三星对4/5nm和8nm新客户提价15\%，并与Anthropic/Meta洽谈2nm AI芯片；英特尔可能获得谷歌TPU订单。但台积电至少领先一个世代，需求持续超过其可提供的领先产能，替代方案对营收影响有限。CoWoS产能2026年底135k WPM→2027年底195k WPM，资本支出2026年560亿美元→2027年680亿美元，毛利率预计从60\%扩至65\%。
{\small\color{refgray}数据：台积电2Q26营收NT\$1.27T，环比+12\%，同比+36\%，比共识高0.5\%\par}
{\small\color{refgray}数据：三星4/5nm和8nm对新客户提价15\%\par}
{\small\color{refgray}数据：CoWoS产能：2026年底135k WPM→2027年底195k WPM\par}
{\small\color{refgray}数据：台积电资本支出：2026年560亿美元→2027年680亿美元\par}
{\small\color{refgray}边际变化：台积电产能瓶颈从市场担忧变为行业公开约束条件，客户主动寻找替代方案，但短期不构成实质性威胁。\par}
\paragraph{MS} 南亚科技所处的利基型DRAM市场，供需紧张格局可能持续到2028年，而非市场预期的2026年见顶。2Q26营收环比+68\%，同比+684\%，EPS NT\$14.66超预期15\%，毛利率从67.9\%跃升至79.5\%。AI基础设施对HBM和DDR5的溢出需求挤压传统DRAM产能，服务器收入占比首次突破20\%。管理层不担心2028年前供应过剩，新增产能主要面向AI定制化产品而非标准型DRAM。
{\small\color{refgray}数据：南亚科2Q26 EPS NT\$14.66，超预期15\%\par}
{\small\color{refgray}数据：毛利率从67.9\%升至79.5\%\par}
{\small\color{refgray}数据：服务器收入占比首次突破20\%\par}
{\small\color{refgray}数据：2027-2029年资本支出160亿美元，80\%用于设备采购\par}
{\small\color{refgray}边际变化：市场通常将大规模资本支出视为周期见顶信号，但新增产能主要面向AI定制化产品，不会直接冲击标准型DRAM供需平衡。\par}
\paragraph{JPM} 南亚科技所处的利基型DRAM市场，供需紧张格局可能持续到2028年，而非市场预期的2026年见顶。2Q26营收环比+68\%，同比+684\%，EPS NT\$14.66超预期15\%，毛利率从67.9\%跃升至79.5\%。AI基础设施对HBM和DDR5的溢出需求挤压传统DRAM产能，服务器收入占比首次突破20\%。管理层不担心2028年前供应过剩，新增产能主要面向AI定制化产品而非标准型DRAM。
{\small\color{refgray}数据：南亚科2Q26 EPS NT\$14.66，超预期15\%\par}
{\small\color{refgray}数据：毛利率从67.9\%升至79.5\%\par}
{\small\color{refgray}数据：服务器收入占比首次突破20\%\par}
{\small\color{refgray}数据：2027-2029年资本支出160亿美元，80\%用于设备采购\par}
{\small\color{refgray}边际变化：市场通常将大规模资本支出视为周期见顶信号，但新增产能主要面向AI定制化产品，不会直接冲击标准型DRAM供需平衡。\par}
\paragraph{MS} 模拟芯片行业正进入更可持续的盈利修复早期阶段。订单改善基于三个关键信号：定价纪律维持、订单取消量有限、客户订单可见度拉长，与投机性补库存行为明显不同。工厂利用率改善、低利用率费用下降，共同指向利润率将进一步扩张。数据中心功率半导体方面，侧挂式电源系统以较低电压独立引入，从3Q26开始出货，有助于对冲800V方案执行不确定性。英飞凌数据中心收入指引FY26 15亿欧元、FY27 25亿欧元。
{\small\color{refgray}数据：英飞凌数据中心收入指引：FY26 15亿欧元、FY27 25亿欧元\par}
{\small\color{refgray}数据：侧挂式电源系统从3Q26开始出货\par}
{\small\color{refgray}数据：订单取消量有限，客户可见度拉长\par}
{\small\color{refgray}边际变化：功率半导体部署不完全绑定于英伟达机架架构演进节奏，存在独立低压出货路径，收入可见度提升。\par}
\paragraph{MS} 市场对台积电3Q26毛利率预期已攀升至69\%-70\%，但海外厂和2nm稀释效应可能使实际值落在67\%-68\%。AI半导体收入五年CAGR从50\%中高段上调至70\%，因记忆体成本上升推动云端资本支出扩张。非AI需求方面，中国智能手机库存修正开始，但影响有限。台积电2027年有5\%-10\%领先制程涨价空间，但2nm初期良率爬坡和海外厂折旧会同步侵蚀利润，毛利率底部约60\%。
{\small\color{refgray}数据：买方预期3Q26毛利率69\%-70\%，MS估算67\%-68\%\par}
{\small\color{refgray}数据：AI半导体收入五年CAGR上调至70\%\par}
{\small\color{refgray}数据：台积电毛利率底部约60\%\par}
{\small\color{refgray}数据：2027年领先制程涨价空间5\%-10\%\par}
{\small\color{refgray}边际变化：市场对涨价传导假设过于线性，低估了海外厂和2nm的稀释效应，但毛利率底部更稳固。\par}
\paragraph{GS} ASMPT的TCB设备受益于AI基建与先进封装趋势，未来三年净利润CAGR可达36\%。全球云服务商资本开支大幅上修：2026-2028年美国CSP上调8\%/27\%/25\%，中国CSP上调37\%/44\%/55\%。AI芯片出货量同期上调14\%/22\%/14\%。TCB设备已获全球IDM 8台重复订单，台积电CoWoS扩产是重要推动力。但当前33倍2027年PE已反映较多预期。
{\small\color{refgray}数据：2026-2028年美国CSP资本开支上调8\%/27\%/25\%\par}
{\small\color{refgray}数据：中国CSP资本开支上调37\%/44\%/55\%\par}
{\small\color{refgray}数据：AI芯片出货量上调14\%/22\%/14\%\par}
{\small\color{refgray}数据：ASMPT净利润CAGR 36\%\par}
{\small\color{refgray}边际变化：中国CSP资本开支上修幅度远超美国，ASMPT增长对国内CSP扩张节奏更为敏感。\par}
\paragraph{GS} 大立光CPO镜头进展明确：最快2027年中进入光纤阵列量产阶段，目前已有客户项目推进。CPO镜头打开精度更高的新市场，光纤阵列对准精度远高于边缘耦合方案，自研主动对准设备构成壁垒。手机镜头业务2026年下半年迎来季节性反弹，毛利率从51\%向56\%以上移动。2027-2028年净利润预测上调1-2\%，基于产品组合升级和CPO早期贡献。
{\small\color{refgray}数据：CPO光纤阵列最快2027年中量产\par}
{\small\color{refgray}数据：毛利率目标从51\%升至56\%以上\par}
{\small\color{refgray}数据：2027-2028年净利润预测上调1-2\%\par}
{\small\color{refgray}数据：2-3周内开始送样，客户认证周期2-4周\par}
{\small\color{refgray}边际变化：CPO镜头从实验室走向量产线，未来18个月产能爬坡可验证，打开新估值空间。\par}
\paragraph{GS} 深南电路AI PCB扩产由客户需求倒逼，2026年资本开支预期上调至60亿元（2025年38亿元），增幅超50\%。公司披露约49亿元定增计划，其中36亿元投入无锡AI PCB产能扩建。2Q26单季净利润中值约13.5亿元，符合预期。增长来自AI服务器、高速交换机、1.6T光模块三大数据中心子领域，产品结构升级推高利润率。
{\small\color{refgray}数据：2026年资本开支预期60亿元，2025年38亿元\par}
{\small\color{refgray}数据：定增49亿元，36亿元用于无锡AI PCB扩产\par}
{\small\color{refgray}数据：2Q26单季净利润中值13.5亿元\par}
{\small\color{refgray}数据：上半年净利润同比增长54\%-69\%\par}
{\small\color{refgray}边际变化：扩产节奏从防御性转为客户需求倒逼的产能前置，资本开支一年内提高近六成。\par}
\paragraph{JPM} 大湾区PCB产业链调研确认，英伟达新一代平台Vera Rubin PCB已进入量产，新供应商加入不自动意味着现有龙头份额流失，最终取决于良率稳定性、材料准备和产能爬坡速度。新架构Kyber NVL144仍在测试，但即使延迟，对2028年AI PCB市场的影响将被现有平台更强需求和ASIC方案加速采用所对冲。1.6T光模块出货增长和CoWoP封装设计推动mSAP工艺需求。
{\small\color{refgray}数据：Vera Rubin PCB已进入量产\par}
{\small\color{refgray}数据：Kyber NVL144架构仍在测试中\par}
{\small\color{refgray}数据：1.6T光模块出货量持续增长\par}
{\small\color{refgray}数据：CoWoP封装推动mSAP工艺需求\par}
{\small\color{refgray}边际变化：竞争逻辑从'谁能进入供应链'切换为'谁能稳定交付'，良率稳定性和交付速度成为关键。\par}
\paragraph{JPM} 亚洲内存板块从6月高点已回撤约30\%，核心矛盾是卖方内存TAM预期跑在了CSP实际资本开支前面。JPM对2026-2027年内存TAM预测为3480-7200亿美元，相当于CSP硬件资本开支的50\%-70\%。约70\%的内存情绪驱动力来自CSP业绩披露前的预期博弈。HBM4E主流方向为12Hi而非16Hi，谷歌TPU和AWS Trainium需求增加。存储级基板/PCB/控制器IC开始紧缺。
{\small\color{refgray}数据：亚洲内存公司从6月高点回撤约30\%\par}
{\small\color{refgray}数据：MSCI亚太指数同期仅跌6\%，费城半导体指数跌11\%\par}
{\small\color{refgray}数据：内存TAM占CSP硬件资本开支50\%-70\%\par}
{\small\color{refgray}数据：HBM4E主流方向为12Hi而非16Hi\par}
{\small\color{refgray}边际变化：市场从'增长加速'阶段进入'盈利耐久性验证'切换期，CSP财报成为短期关键催化剂。\par}
\paragraph{MS} 博通在谷歌TPU供应链中仍将维持约80\%份额，联发科参与不构成结构性取代。市场低估了在规模化TPU平台上替换博通的难度：包括HBM供应锁定、封装执行经验、系统级认证和大规模部署稳定性。博通已签好长期HBM合同锁定低价，联发科在公开市场采购HBM成本更高。2nm TPU封装路径（如EMIB）仍不确定，CoWoS仍是更可靠方案。
{\small\color{refgray}数据：博通预计维持约80\% TPU份额\par}
{\small\color{refgray}数据：联发科在3nm TPU上有可信机会\par}
{\small\color{refgray}数据：博通已锁定长期低价HBM合同\par}
{\small\color{refgray}数据：2nm TPU封装路径仍不确定\par}
{\small\color{refgray}边际变化：AI加速器竞争从'谁能做出芯片'转向'谁能大规模可靠部署芯片'，博通的供应链关系和封装经验构成护城河。\par}
\subsubsection*{关键数据}
\begin{itemize}
  \item 台积电2Q26营收NT\$1.27T，环比+12\%，同比+36\%，比共识高0.5\%
  \item 南亚科2Q26 EPS NT\$14.66，超预期15\%，毛利率从67.9\%升至79.5\%
  \item 亚洲内存公司从6月高点回撤约30\%，MSCI亚太指数同期仅跌6\%
  \item 2026-2028年中国CSP资本开支上调37\%/44\%/55\%，远超美国CSP的8\%/27\%/25\%
  \item 博通预计维持约80\%谷歌TPU份额，联发科参与不构成结构性取代
  \item 台积电CoWoS产能2026年底135k WPM→2027年底195k WPM
  \item 模拟芯片订单改善基于定价纪律维持、订单取消量有限、客户可见度拉长
  \item 英飞凌数据中心收入指引FY26 15亿欧元、FY27 25亿欧元
\end{itemize}
\begin{figure}[htbp]
\centering
\includegraphics[width=0.88\linewidth]{figures/fig_001.jpg}
\caption{EXHIBIT 1 - EXHIBIT 1: TSMC reported 2Q26 revenue at NT\textbackslash{}\$1.27T, reaching its guidance high-end \& exceeding consensus by c. 0.5\%. Our analysis also suggests FX has no bearing on this finding. TSMC 2Q26 Revenue}
\end{figure}
\begin{figure}[htbp]
\centering
\includegraphics[width=0.88\linewidth]{figures/fig_171.jpg}
\caption{Figure 2 - Figure 2: Global memory makers' share price performance including SOX (Philadelphia Semiconductor index) \#\# Q\&A highlights}
\end{figure}
\begin{figure}[htbp]
\centering
\includegraphics[width=0.88\linewidth]{figures/fig_173.jpg}
\caption{Figure 1 - Figure 1: Global memory makers' share price performance including SOX (Philadelphia Semiconductor Index) \%}
\end{figure}
\begin{figure}[htbp]
\centering
\includegraphics[width=0.88\linewidth]{figures/fig_187.jpg}
\caption{Exhibit 1 - Exhibit 1: AVGO and NVDA performance have lagged other AI Semis \#\# Broadcom Inc. (AVGO.O, AVGO US)}
\end{figure}
{\small\color{refgray}本节综合 11 篇研究；涉及 Bernstein、MS、GS、JPM。}
\newpage
\subsection{AI基础设施与平台分化：从算力堆叠到生态变现}
\textbf{AI基础设施投资正从概念验证转向规模部署，但受益结构出现分化：算力层（GPU/交换机）需求强劲，但平台层（Naver、SAP）和内容层（B站、短剧）的差异化价值开始被重新定价。市场共识是AI资本开支周期加速，分歧在于谁能在生态和系统优化中捕获超额利润。}
\subsubsection*{主流共识}
\begin{itemize}
  \item AI基础设施资本开支周期加速，中国云厂商AI capex上调37\%-55\%（GS R054），全球AI服务器芯片预期上调14\%-22\%。
  \item Scale-Up网络市场规模因集群规模膨胀而暴增，MS将2030年预估从170亿美元上调至730亿美元（R020）。
  \item 英伟达短期业绩强劲（Q2增长95\%），但增量买家不足导致股价相对落后（MS R015）。
  \item AI内容（短剧）成本优势显著（为真人1/3-1/4），但平台主动控制占比以保护广告价值（NOM R022）。
\end{itemize}
\subsubsection*{分歧与条件差异}
\begin{itemize}
  \item Naver的AI工厂应参照CoreWeave（算力租赁）还是Nebius（全栈平台）估值？Bernstein认为后者更合理（R007）。
  \item 铜缆与光模块在Scale-Up网络中的替代节奏：MS认为铜缆寿命被低估，CPO规模化要到2029年后（R020）；市场对CPO延迟的担忧可能过度。
  \item 开源模型是否拉平竞争？NOM专家指出DeepSeek的成本优势来自系统级优化（缓存、延迟、硬件利用率），无法通过开源复制（R023）。
  \item 苹果iPhone 18是否提价200美元？MS基于弹性分析（0.2-0.5）认为可行且利润上行（R062）；市场预期100-150美元。
\end{itemize}
\subsubsection*{机构观点}
\paragraph{Bernstein} Naver的AI工厂在基础设施层面与CoreWeave相似（SPV融资、GPU原生架构），但战略意图指向Nebius——构建区域AI全栈平台，整合搜索、韩语AI、企业服务和主权AI需求。长期估值不应参照GPU租赁公司，而应参照从基础设施走向平台的公司。Naver已锁定初始200 MW电力容量，目标5-6年内新增20万亿韩元收入。
{\small\color{refgray}数据：55 MW起步规模，后续扩张\par}
{\small\color{refgray}数据：目标5-6年新增20万亿韩元收入\par}
{\small\color{refgray}数据：初始200 MW电力容量已锁定\par}
{\small\color{refgray}边际变化：市场此前将Naver的neocloud与CoreWeave类比，Bernstein修正为平台型估值逻辑，强调区域AI栈的差异化。\par}
\paragraph{MS} 海光信息（Hygon）的“CPU+GPU”一体化模式与纯GPU公司形成差异化，可覆盖通用服务器和AI训练/推理。Sugon 8000十万卡集群全部采用国产芯片，验证了Hygon生态的集成能力。市场可能低估CPU在AI基础设施中的结构性价值，尤其是在中国语境下。
{\small\color{refgray}数据：Sugon 8000：国内首个10万卡AI超算集群\par}
{\small\color{refgray}数据：六类核心芯片全部自研：CPU、GPGPU、交换芯片、网卡\par}
{\small\color{refgray}数据：Hygon第一大股东Sugon持股27.96\%\par}
{\small\color{refgray}边际变化：从单一GPU对标转向平台型竞争逻辑，强调CPU在AI部署中的结构性价值。\par}
\paragraph{MS} 英伟达面临的最大挑战不是需求或竞争，而是庞大市值导致增量买家不足。管理层正接触价值型和收入型账户，当前季度仍保持95\%同比增长，FCF回报率超5\%。AI数据中心地理回迁趋势加速，各国为本国企业预留资源。内存短缺将持续数年，NVDA尝试“用更少内存做更多事”。
{\small\color{refgray}数据：Q2同比增长95\%\par}
{\small\color{refgray}数据：FCF回报率超5\%\par}
{\small\color{refgray}数据：目标价\$288 vs 当前\$211\par}
{\small\color{refgray}边际变化：从纯成长股叙事转向价值+成长双重属性，强调增量买家不足是股价相对落后的主因。\par}
\paragraph{MS} Scale-Up网络市场规模从170亿美元（2029年）上调至730亿美元（2030年），核心驱动力是集群规模膨胀（从8 GPU到1000+ GPU）。铜缆在短距离互联中的寿命被低估，CPO规模化部署要到2029年后。近期市场对CPO延迟的担忧过度放大了短期不确定性。
{\small\color{refgray}数据：2030年Scale-Up市场预估730亿美元（同比4倍）\par}
{\small\color{refgray}数据：NVL72：72 GPU/机柜；Rubin Ultra：576 GPU；Feynman：1000+ GPU\par}
{\small\color{refgray}数据：每GPU光学引擎数从2个（NVL72）增至70个（Feynman全CPO）\par}
{\small\color{refgray}边际变化：从线性需求外推修正为架构性变化驱动，强调集群规模膨胀是TAM暴增的真正推手。\par}
\paragraph{NOM} 红果短剧平台主动将AI内容占比控制在30\%以下，因AI短剧吸引率仅2-8\%（优质真人10-20\%），流行周期2-4周，侵蚀广告价值。平台通过推荐权重、分成调整和审核将70\%流量分配给真人内容。AI短剧市场规模预计2026年达267亿元（同比+80\%），但平台干预可能影响AI工具变现路径。
{\small\color{refgray}数据：红果DAU 1.2亿，MAU 3.3亿\par}
{\small\color{refgray}数据：AI内容占比<30\%，目标控制在此水平\par}
{\small\color{refgray}数据：AI短剧吸引率2-8\% vs 真人10-20\%\par}
{\small\color{refgray}数据：AI短剧流行周期2-4周\par}
{\small\color{refgray}边际变化：从AI内容增长叙事转向平台主动管理内容生态以保护广告价值，强调成本优势不等于商业优势。\par}
\paragraph{NOM} 大模型价格分化：基础模型降价引流，高端模型涨价锁定企业客户。DeepSeek的成本优势来自系统级优化（缓存效率、延迟、硬件利用率），无法通过开源复制。芯片供应紧张推动国产硬件在推理场景性价比改善。金融和办公效率是商业化潜力最大的行业。
{\small\color{refgray}数据：高端英伟达芯片报价过去一年大幅上涨\par}
{\small\color{refgray}数据：国产加速器涨幅相对温和\par}
{\small\color{refgray}数据：DeepSeek推理成本显著低于同类模型\par}
{\small\color{refgray}数据：企业客户要求12-18个月AI投入回报\par}
{\small\color{refgray}边际变化：从开源拉平竞争叙事转向系统优化护城河，强调芯片供应紧张重塑国产硬件性价比。\par}
\paragraph{NOM} 腾讯2Q26F收入预计增长10\%至2030亿，非IFRS净利润增长6\%至670亿，略低于共识3\%，因AI投入压制利润率。混元3.0正式版获开发者认可，WorkBuddy月访问量885万。广告（+18\%）和云（+22\%）是短期增长引擎，但AI投入导致利润率回落0.4个百分点至37.1\%。
{\small\color{refgray}数据：2Q26F收入2030亿（+10\% y-y）\par}
{\small\color{refgray}数据：非IFRS净利润670亿（+6\% y-y）\par}
{\small\color{refgray}数据：非IFRS经营利润率37.1\%（低于共识37.4\%）\par}
{\small\color{refgray}数据：广告收入+18\%，云收入+22\%\par}
{\small\color{refgray}边际变化：从AI故事验证转向利润表影响显现，强调AI投入对短期利润率的压制。\par}
\paragraph{NOM} 紫光股份1H26归母净利润同比增长83.5\%-122.9\%，核心驱动力来自中国AIDC资本开支加速和国产替代。网络设备FY26F增长27\%，服务器增长40\%，高端产品升级可抵消白牌对利润率的稀释。海外市场突破（新华三整合）打开新空间。
{\small\color{refgray}数据：1H26归母净利润19.1-23.2亿元（+83.5\%-122.9\%）\par}
{\small\color{refgray}数据：经常性净利润+44\%-81\%\par}
{\small\color{refgray}数据：网络设备FY26F +27\%，服务器+40\%\par}
{\small\color{refgray}数据：FY26-28F网络设备/服务器CAGR 15\%/25\%\par}
{\small\color{refgray}边际变化：从服务器出货量关注转向网络交换机技术升级和海外突破，强调高端产品定价权。\par}
\paragraph{Bernstein} SAP最核心增长驱动力是S/4HANA迁移，而非AI产品变现。ECC安装基数仅55\%完成迁移，剩余45\%客户构成可预见云收入转化机会至2027年。CCB增速稳定在24-25\%将巩固市场信心。AI是加速器而非独立收入故事。
{\small\color{refgray}数据：ECC迁移率55\%\par}
{\small\color{refgray}数据：CCB增速1Q26 +25\%，预期2Q26 +24\%\par}
{\small\color{refgray}数据：云收入增长可支撑至2027年\par}
{\small\color{refgray}边际变化：从AI叙事主导转向S/4迁移结构性周期，强调ERP系统替换的刚性需求。\par}
\paragraph{GS} AI基础设施资本开支周期加速，中国云厂商AI capex上调37\%-55\%。圣邦股份（模拟芯片）向AI数据中心延伸，200+料号进入数据中心，2026-2028年净利润上调1\%/5\%/9\%。锐捷网络（白牌交换机）受益于云厂商超节点建设，但估值已回均值，下调至中性。
{\small\color{refgray}数据：全球AI服务器芯片预期上调14\%-22\%\par}
{\small\color{refgray}数据：中国云厂商AI capex上调37\%-55\%\par}
{\small\color{refgray}数据：圣邦股份200+料号进入数据中心\par}
{\small\color{refgray}数据：网络/计算模拟IC市场2030年达1330亿元（CAGR 16\%）\par}
{\small\color{refgray}边际变化：从边际修正到周期斜率重估，强调中国云厂商投入力度超预期。\par}
\paragraph{MS} 苹果提价策略从追求收入转向保护毛利率。iPhone需求弹性仅0.2-0.5（远低于Mac 0.8和iPad 1.0），iPhone 18 Pro可能提价200美元。Mac/iPad提价15-30\%主要对冲存储芯片成本上升（DRAM +190\%，NAND +180\%）。
{\small\color{refgray}数据：iPhone需求弹性0.2-0.5\par}
{\small\color{refgray}数据：Mac/iPad提价15-30\%\par}
{\small\color{refgray}数据：FY27 DRAM成本+190\%，NAND+180\%\par}
{\small\color{refgray}数据：iPhone 18 Pro可能提价200美元\par}
{\small\color{refgray}边际变化：从收入增长叙事转向毛利率保护，强调iPhone需求刚性被低估。\par}
\paragraph{MS} B站从内容分发平台向内容社区完成身份确认。800万用户关注同一创作者超10年，65\%观看时长在5分钟以上视频，月均互动170亿次。AI生成广告封面提升点击率8ppt，全面推广后eCPM可提升5\%。用户平均年龄27岁，55\%在一二线城市。
{\small\color{refgray}数据：800万用户关注同一创作者>10年\par}
{\small\color{refgray}数据：65\%观看时长在5分钟以上视频\par}
{\small\color{refgray}数据：月均互动170亿次\par}
{\small\color{refgray}数据：用户平均年龄27岁，55\%在一二线城市\par}
{\small\color{refgray}边际变化：从DAU增长叙事转向社区粘性和长生命周期内容价值，强调商业化飞轮。\par}
\subsubsection*{关键数据}
\begin{itemize}
  \item Scale-Up网络2030年市场预估730亿美元（MS R020）
  \item 中国云厂商AI capex上调37\%-55\%（GS R054）
  \item 英伟达Q2同比增长95\%，FCF回报率超5\%（MS R015）
  \item 红果短剧AI内容占比<30\%，AI短剧吸引率2-8\% vs 真人10-20\%（NOM R022）
  \item SAP ECC迁移率仅55\%，支撑云收入至2027年（Bernstein R029）
  \item 苹果iPhone需求弹性0.2-0.5，iPhone 18 Pro可能提价200美元（MS R062）
  \item B站800万用户关注同一创作者超10年（MS R064）
  \item 紫光股份1H26归母净利润+83.5\%-122.9\%（NOM R028）
  \item 腾讯2Q26F非IFRS经营利润率37.1\%，低于共识（NOM R026）
  \item Naver AI工厂目标5-6年新增20万亿韩元收入（Bernstein R007）
\end{itemize}
\begin{figure}[htbp]
\centering
\includegraphics[width=0.88\linewidth]{figures/fig_076.jpg}
\caption{Exhibit 1 - Exhibit 1: Scale-Up Market Estimate Has 4x'd Y/Y Exhibit 2: CPO Penetration to Come Later to Scale-Up, But Post 2028, Should Be Meaningful Adoption \#\# Analysis}
\end{figure}
\begin{figure}[htbp]
\centering
\includegraphics[width=0.88\linewidth]{figures/fig_156.jpg}
\caption{Exhibit 1 - Exhibit 1: China CSP capex trend Data include Bytedance, Tencent, Alibaba, Baidu. We calculate Bytedance (private) relevant operating metrics by analyzing the industry and companies that our China Internet team cover, and then ex}
\end{figure}
\begin{figure}[htbp]
\centering
\includegraphics[width=0.88\linewidth]{figures/fig_087.jpg}
\caption{EXHIBIT 4 - EXHIBIT 4: SAP: total S/4 customers added EXHIBIT 5: SAP: ECC clients to be migrated to S/4 EXHIBIT 6: SAP: new customers having acquired S/4}
\end{figure}
\begin{figure}[htbp]
\centering
\includegraphics[width=0.88\linewidth]{figures/fig_175.jpg}
\caption{Exhibit 1 - Exhibit 3: Despite the limited sample size, our Sector Quant team's work suggests that iPhones (0.2-0.4) and Macs (0.8) have historically exhibited relatively inelastic demand, while iPads (1.0) have shown roughly unitary price ela}
\end{figure}
{\small\color{refgray}本节综合 12 篇研究；涉及 Bernstein、MS、NOM、GS。}
\newpage
\subsection{科技与互联网（2）：百度AI转型的结构性错位}
\textbf{百度正经历广告收入持续探底与AI云高速增长的结构性错位，AI研究投入加大进一步压缩核心经营利润，短期业绩承压但长期AI商业化路径逐步清晰。}
\subsubsection*{主流共识}
\begin{itemize}
  \item 百度广告业务受宏观环境和AI转型双重影响，二季度核心广告收入同比下滑约18.5\%，复苏节奏缓慢。
  \item AI云基础设施收入保持高速增长，GPU云服务维持三位数增速，但利润贡献仍以IaaS层为主。
  \item 下半年AI研究投入加大，将压缩核心经营利润，全年核心Non-GAAP营业利润预计同比下降3-4\%。
\end{itemize}
\subsubsection*{分歧与条件差异}
\begin{itemize}
  \item AI广告产品（如数字人、智能体）已贡献18\%广告收入，但能否快速替代传统搜索广告下滑存在分歧。
  \item AI云收入增长强劲，但利润转化时间不确定，部分观点认为IaaS层利润改善需更长时间。
\end{itemize}
\subsubsection*{机构观点}
\paragraph{MS} MS预计百度二季度核心收入同比下滑4\%，广告收入同比降18.5\%（较一季度-21\%略有收窄），但下半年复苏空间有限。AI云基础设施收入76.3亿元，同比增56\%，GPU云服务三位数增长，但利润转化仍以IaaS为主。下半年AI研究投入加大，核心经营利润率预计降至14.7\%（同比降2个百分点），全年核心Non-GAAP营业利润同比降3-4\%。非核心资产剥离推进中，年底有望首次派息，双重主要上市或在2027年一季度。
{\small\color{refgray}数据：核心广告收入130亿元，同比-18.5\%\par}
{\small\color{refgray}数据：AI云基础设施收入76.3亿元，同比+56\%\par}
{\small\color{refgray}数据：GPU云服务三位数增长\par}
{\small\color{refgray}数据：核心经营利润率14.7\%，同比-2个百分点\par}
{\small\color{refgray}边际变化：相比此前预期，MS下调了广告收入复苏节奏，并强调下半年AI研究投入加大将进一步压缩利润，同时首次提及派息和双重主要上市时间表。\par}
\subsubsection*{关键数据}
\begin{itemize}
  \item 核心广告收入130亿元，同比-18.5\%
  \item AI云基础设施收入76.3亿元，同比+56\%
  \item GPU云服务三位数增长
  \item 核心经营利润率14.7\%，同比-2个百分点
  \item 全年核心Non-GAAP营业利润同比-3\textasciitilde{}4\%
  \item AI广告产品贡献核心广告收入18\%
\end{itemize}
{\small\color{refgray}本节综合 1 篇研究；涉及 MS。}
\newpage
\subsection{汽车与零部件：结构性分化与定价权验证}
\textbf{全球汽车与零部件行业正经历结构性分化：欧洲传统车企受中国需求下滑与成本压力挤压利润率，而中国新能源品牌凭借产品换代与品牌溢价实现逆势增长；轮胎行业短期销量承压，但成本冲击恰是验证定价权的时刻。}
\subsubsection*{主流共识}
\begin{itemize}
  \item 中国新能源市场内部格局分化，比亚迪份额承压，新势力与国企反攻，高端纯电市场集中度提升。
  \item 欧洲传统车企（宝马、奔驰）受中国销量下滑拖累，利润率承压，产能东迁成为应对成本压力的共同策略。
  \item 特斯拉汽车业务回暖但非估值核心，市场焦点转向Robotaxi和Optimus的物理AI进展。
\end{itemize}
\subsubsection*{分歧与条件差异}
\begin{itemize}
  \item 燃油车渗透率反弹是供给端主动去库存的结果还是需求趋势逆转？Citi认为库存去化速度（燃油车库存月数从5.0降至4.0）表明供给端调整，而非需求系统性回暖。
  \item 日本零部件供应商的定价权是否正在被中国本土供应链替代？MS指出日系整车厂加速采用中国零部件，日本供应商溢价空间收窄。
  \item 轮胎行业成本冲击是短期节奏问题还是结构性影响？Bernstein认为历史经验显示定价权可应对，预计2027上半年完全消化成本。
\end{itemize}
\subsubsection*{机构观点}
\paragraph{Bernstein} 瑞安航空通过机队升级（MAX-10单位成本较NG低14\%）、发动机维修自营（IRR 23\%）和机场议价权，将成本优势锁定为长期结构性壁垒，与其他航司成本通胀不确定性形成鲜明对比。全球轮胎行业二季度销量全面转负（北美乘用车替换市场同比-4\%），但成本冲击恰是验证定价权的时刻，参考2021-2022年成本上升期轮胎公司表现优于市场并触及估值高点。
{\small\color{refgray}数据：MAX-10单位成本较737-NG低14\%，燃油成本下降27\%，维护成本下降17\%\par}
{\small\color{refgray}数据：发动机维修自营项目IRR 23\%，累计节省5-8亿欧元\par}
{\small\color{refgray}数据：北美乘用车替换市场Q2同比-4\%，中国替换市场+13\%\par}
{\small\color{refgray}数据：北美轮胎进口量同比降幅达两位数，终端销售仅降4\%\par}
{\small\color{refgray}边际变化：从静态成本优势转向可量化的长期结构性壁垒锁定；轮胎行业从关注成本压力转向验证定价权。\par}
\paragraph{MS} 日本车企在华销量加速下滑（丰田-27\%、本田-47\%、日产-32\%），零部件供应商正失去订单量、定价权和供应链位置；日系整车厂加速采用中国本土零部件，日本供应商溢价空间收窄。奔驰匈牙利工厂成本比德国低70\%，产能东迁策略将全球产能从250万辆调整至200-220万辆，以应对德国本土成本竞争力下降。
{\small\color{refgray}数据：丰田Q2在华销量同比-27\%，本田-47\%，日产-32\%\par}
{\small\color{refgray}数据：三家合计Q2销量从2024年同期约79万辆降至约51万辆\par}
{\small\color{refgray}数据：奔驰匈牙利工厂成本比德国低70\%，产能弹性从20万辆提升至40万辆\par}
{\small\color{refgray}数据：奔驰全球产能目标从250万辆降至200-220万辆\par}
{\small\color{refgray}边际变化：日本零部件供应商从受益于日元贬值转向中国业务量下降侵蚀利润；奔驰从被动应对转向主动东迁以锁定成本优势。\par}
\paragraph{MS} 特斯拉Q2交付48万辆超预期18\%，但汽车业务已非估值核心，市场焦点转向Robotaxi和Optimus的物理AI进展。FSD累计里程突破100亿英里，但消费者版无人FSD最早4Q26。中国电动车6月零售定价整体稳定，比亚迪综合价格环比仅降0.4\%，蔚来因ES9上市环比涨3.5\%，新车型推出对冲了老款促销压力。
{\small\color{refgray}数据：特斯拉Q2交付48万辆，高于市场预期40.6万辆\par}
{\small\color{refgray}数据：FSD累计行驶里程突破100亿英里，每天行驶约2670万英里\par}
{\small\color{refgray}数据：比亚迪6月综合零售价格环比-0.4\%，汉系列折扣达20\%\par}
{\small\color{refgray}数据：蔚来6月综合价格环比+3.5\%，ES9上市拉高入门价格\par}
{\small\color{refgray}边际变化：特斯拉汽车业务从增长引擎转为安全垫；中国电动车市场从全面降价担忧转向产品结构支撑定价。\par}
\paragraph{Citi} 6月中国燃油车渗透率环比回升2个百分点至39.7\%，但库存去化速度更快（燃油车库存月数从5.0降至4.0），回升更多是供给端调整而非需求趋势逆转。新能源内部格局分化，比亚迪市占率22.1\%同比降8.3个百分点，蔚来、小鹏、上汽份额环比提升。
{\small\color{refgray}数据：6月燃油车保险零售渗透率39.7\%，环比+2ppt\par}
{\small\color{refgray}数据：乘用车整体库存降至2.7个月，环比-0.5个月\par}
{\small\color{refgray}数据：燃油车库存月数从5.0降至4.0，新能源从2.1降至1.9\par}
{\small\color{refgray}数据：比亚迪新能源市占率22.1\%，同比-8.3ppt\par}
{\small\color{refgray}边际变化：从关注新能源渗透率持续上升转向燃油车库存去化速度超预期；从比亚迪一家独大转向新势力与国企份额反攻。\par}
\paragraph{DB} 宝马Q2汽车业务利润率预计降至1.2\%（去年同期5.4\%），中国销量同比降30\%是最大单一变量，抵消了欧洲和美洲的增长。定价环境恶化、折旧摊销和原材料成本进一步压缩利润空间，IEEPA退税（低三位数百万欧元）远不足以弥补下滑。
{\small\color{refgray}数据：宝马Q2汽车业务利润率预计1.2\%，同比-425bps\par}
{\small\color{refgray}数据：Q2全球零售销量59.1万辆，同比-5\%，中国-30\%\par}
{\small\color{refgray}数据：汽车业务营业利润预计3.32亿欧元，同比-79\%\par}
{\small\color{refgray}数据：集团营业利润预计12.47亿欧元，同比-53\%\par}
{\small\color{refgray}边际变化：从关注欧洲传统车企的电动化转型转向中国利润池从增量变为减量；从季度波动预期转向结构性定价权变化。\par}
\paragraph{GS} 蔚来在行业零售同比下滑14\%背景下交付19.1万辆同比+67\%，40万元以上纯电市场份额从6\%升至39\%。产品换代（ES8月销超1万辆）和品牌溢价（净推荐值行业第一）推动非GAAP盈利有望在2026年转正，自由现金流从-31亿转正至+121亿。
{\small\color{refgray}数据：蔚来1H26交付19.1万辆，同比+67\%\par}
{\small\color{refgray}数据：40万元以上纯电市场份额从1H25的6\%升至1H26的39\%\par}
{\small\color{refgray}数据：ES8月均销量超1万辆，在40万元以上SUV细分市场排名第一\par}
{\small\color{refgray}数据：2026年预计Non-GAAP盈利16亿（2025年亏损124亿）\par}
{\small\color{refgray}边际变化：从亏损换销量转向产品换代驱动盈利转正；从市场折价转向基本面改善与估值脱节修复。\par}
\paragraph{JEF} 特斯拉Q2交付48万辆超预期，但产品创新节奏放缓，汽车商品化趋势加速。Cybercab产量仅8822辆，Robotaxi规模化非短期可实现。Q2 EBIT预期上调至14.5亿美元（利润率5.1\%），但远期盈利仍低于市场共识，因Robotaxi和Optimus早期影响利润。
{\small\color{refgray}数据：特斯拉Q2交付480,128辆，高于市场预期40.6万辆\par}
{\small\color{refgray}数据：Cybercab Q2产量8,822辆，交付12,364辆\par}
{\small\color{refgray}数据：Q2 EBIT预期14.5亿美元，利润率5.1\%\par}
{\small\color{refgray}数据：2026年EBIT预期上调4\%至62亿\par}
{\small\color{refgray}边际变化：从关注交付量超预期转向产品创新节奏放缓与Robotaxi规模化时间线；从汽车业务估值转向物理AI进展决定长期价值。\par}
\subsubsection*{关键数据}
\begin{itemize}
  \item 宝马Q2中国销量同比-30\%，汽车业务利润率降至1.2\%
  \item 蔚来1H26交付19.1万辆同比+67\%，40万元以上纯电份额39\%
  \item 特斯拉Q2交付48万辆超预期18\%，FSD累计里程突破100亿英里
  \item 中国6月燃油车渗透率39.7\%环比+2ppt，库存月数从5.0降至4.0
  \item 奔驰匈牙利工厂成本比德国低70\%，产能弹性至40万辆
  \item 瑞安航空MAX-10单位成本较NG低14\%，发动机维修自营IRR 23\%
  \item 比亚迪6月综合价格环比-0.4\%，蔚来因ES9上市环比+3.5\%
  \item 日本三家车企Q2在华销量合计约51万辆，同比降幅超30\%
\end{itemize}
\begin{figure}[htbp]
\centering
\includegraphics[width=0.88\linewidth]{figures/fig_045.jpg}
\caption{Exhibit 1 - Exhibit 1: New vehicles sales of 3 Japanese OEMs in China (monthly) Exhibit 2: New vehicle sales of Japanese OEMs in China (quarterly) MS MUFG SECURITIES CO., LTD.+}
\end{figure}
\begin{figure}[htbp]
\centering
\includegraphics[width=0.88\linewidth]{figures/fig_056.jpg}
\caption{Exhibit 1 - Exhibit 1: MBG new model releases by cohort: FY25-26 models to drive volumes Exhibit 2: New models could represent a higher proportion of sales than at peers Exhibit 3: MBG Guidance in perspective}
\end{figure}
\begin{figure}[htbp]
\centering
\includegraphics[width=0.88\linewidth]{figures/fig_081.jpg}
\caption{Exhibit 2 - Exhibit 2: Modestly Increasing Auto Delivery Forecasts \#\# 2) Pulse Check on Autonomy. TSLA FSD hit a key autonomy milestone in 2Q; can the momentum continue? We look to an update on Unsupervised FSD, v15 rollout, and global appro}
\end{figure}
\begin{figure}[htbp]
\centering
\includegraphics[width=0.88\linewidth]{figures/fig_095.jpg}
\caption{Figure 1 - Figure 1. NEV Insurance Retail by OEM © 2026 Citi Inc. No redistribution without Citi's written permission. Figure 2. Sector Volume and Penetration by Fuel Type (Domestically Produced) © 2026 Citi Inc. No redistribution without Citi's written permission. Figur}
\end{figure}
{\small\color{refgray}本节综合 10 篇研究；涉及 Bernstein、MS、Citi、DB、GS、JEF。}
\newpage
\subsection{消费分化加剧：美妆韧性、奢侈品预期触底、中国结构性机会}
\textbf{全球消费市场呈现显著分化：西欧美妆个护销量驱动增长，奢侈品板块预期已充分定价疲软，中国宠物食品与零食行业出现结构性改善信号，日本便利店网络价值重估。}
\subsubsection*{主流共识}
\begin{itemize}
  \item 西欧消费品市场整体增速放缓，食品、家庭护理和消费者健康品类价值增速均出现回落，消费者趋于理性。
  \item 中国国内奢侈品消费仍在调整，但韩国和美国市场正在接棒，成为增长新引擎。
  \item 欧洲奢侈品板块估值接近历史低位，投资者情绪疲软，但一旦出现边际改善，资金可能轮动回归。
\end{itemize}
\subsubsection*{分歧与条件差异}
\begin{itemize}
  \item 美妆个护板块逆势加速，价值增速从3.6\%提升至4.1\%，且完全由销量驱动（+4.7\%），而非涨价，与其他品类形成鲜明对比。
  \item 中国宠物食品线上渠道分化：抖音增长34\% vs 天猫淘宝仅增3\%，本土品牌内部分化加剧，非上市品牌增速（30\%）远超全球品牌（2\%）。
  \item 京东即时零售亏损大幅收窄（从135亿降至60亿），盈利韧性被低估，但市场对其收入下滑（-5.5\%）仍存担忧。
  \item 日本便利店业态虽面临客流下降，但网络衍生价值（支付、零售媒体）被资本重估，软银等巨头持续投入。
\end{itemize}
\subsubsection*{机构观点}
\paragraph{BARC} 欧洲奢侈品2026年第二季度有机增长预计加速至4.7\%，略高于市场预期的4.3\%，但加速来源并非中国，而是美国（信用卡消费增速从6\%升至16\%）和韩国（百货奢侈品销售同比+37\%）。行业内部分化加剧，历峰珠宝预计增长16\%，而Gucci可能下降5\%。中东地缘风险是最大增量不确定性，但5月中东飞往欧洲航班量回升，边际改善信号出现。
{\small\color{refgray}数据：覆盖公司有机增长4.7\%，高于市场预期4.3\%\par}
{\small\color{refgray}数据：美国信用卡消费增速从Q1的6\%升至Q2的16\%\par}
{\small\color{refgray}数据：韩国百货奢侈品销售5月同比+37\%\par}
{\small\color{refgray}数据：中国5月零售总额同比-0.6\%，两年来首次负增长\par}
{\small\color{refgray}边际变化：增长引擎从中国转向美国和韩国，中东风险边际改善（航班量回升），但中国国内消费仍在调整。\par}
\paragraph{Bernstein} 西欧消费品市场分化明显：美妆个护逆势加速，价值增速从3.6\%升至4.1\%，完全由销量驱动（+4.7\%），而食品价值增速从4.2\%骤降至1.9\%，消费者健康从+5.0\%跌至-1.3\%。日本便利店行业正成为物理AI规模化部署的关键场景，软银等巨头投入约1.4万亿日元，看重的是网络衍生价值（Seven Bank ATM每平方米利润是门店的40倍，零售媒体市场预计2029年达1.3万亿日元），而非零售本身。
{\small\color{refgray}数据：美妆个护价值增速+4.1\%，销量增速+4.7\%\par}
{\small\color{refgray}数据：食品价值增速从+4.2\%降至+1.9\%\par}
{\small\color{refgray}数据：消费者健康从+5.0\%降至-1.3\%\par}
{\small\color{refgray}数据：日本便利店客流量指数从100降至87.5\par}
{\small\color{refgray}边际变化：美妆个护韧性由销量驱动而非涨价，日本便利店价值从零售转向网络衍生层（支付、媒体、AI测试床）。\par}
\paragraph{NOM} 京东2026年第二季度盈利韧性超预期：零售收入预计仅下降5.5\%（好于市场预期的-7.5\%），经营利润率稳定在4.5\%，经营利润132亿元比彭博一致预期高3\%。即时零售亏损从去年同期的135亿元大幅收窄至约60亿元，是利润改善的最大变量。电子品类收入下滑13\%，但超市业务增长9\%，在一般商品品类中占比接近45\%，部分对冲了整体下滑。
{\small\color{refgray}数据：京东零售收入预计同比-5.5\%，市场预期-7.5\%\par}
{\small\color{refgray}数据：经营利润率4.5\%，经营利润132亿元\par}
{\small\color{refgray}数据：即时零售亏损从135亿收窄至60亿\par}
{\small\color{refgray}数据：电子品类收入-13\%，超市业务+9\%\par}
{\small\color{refgray}边际变化：即时零售亏损以季度为单位快速收窄，盈利韧性被低估；超市业务成为增长核心引擎。\par}
\paragraph{DB} 欧洲奢侈品板块二季度业绩预期已充分定价疲软环境，投资者情绪接近历史底部。珠宝是相对亮点，历峰珠宝部门预计增长13-15\%，但市场隐含预期可能已跑在前面（历峰交易在约25倍2027年预期市盈率）。中国仍是最大变量，4-5月进口数据继续走弱，美国保持强劲但未加速。板块估值接近历史低位，一旦出现环比改善，资金可能轮动回归，LVMH是首选复苏标的。
{\small\color{refgray}数据：历峰珠宝部门市场隐含预期增长+15\%或更高\par}
{\small\color{refgray}数据：历峰交易在约25倍2027年预期市盈率\par}
{\small\color{refgray}数据：中国4-5月奢侈品进口数据继续走弱\par}
{\small\color{refgray}数据：美国保持强劲但未加速\par}
{\small\color{refgray}边际变化：市场预期已足够保守，珠宝预期可能过高；中国未见实质性改善，但估值底部提供安全边际。\par}
\paragraph{GS} 洽洽食品2026年第二季度净利润中位数8400万元，超预期12\%，驱动因素为竞争格局改善和葵花籽采购成本有利。利润率改善幅度（隐含净利润率6.4\% vs 预期5.7\%）大于收入，说明成本端和费用端贡献更直接。中国宠物食品线上渠道分化加剧：二季度抖音GMV增长34\%，天猫淘宝仅增3\%；本土品牌内部分化，非上市品牌增长30\%，全球品牌仅增2\%，头部集中度提升。
{\small\color{refgray}数据：洽洽Q2净利润中位数8400万，超预期12\%\par}
{\small\color{refgray}数据：隐含净利润率6.4\%，高于预期的5.7\%\par}
{\small\color{refgray}数据：宠物食品线上总GMV同比+14\%\par}
{\small\color{refgray}数据：抖音GMV+34\%，天猫淘宝+3\%\par}
{\small\color{refgray}边际变化：洽洽利润改善来自竞争趋缓和成本红利，可持续性强；宠物食品渠道红利向头部集中，抖音不再是低门槛渠道。\par}
\paragraph{JEF} Shopify增长逻辑正从拓展商户数量转向挖掘现有商户价值：合作伙伴激励改革（前四年高佣金、之后归零）短期推高获客，长期优化成本；定价调整窗口打开，若2027年对非Plus套餐执行类似2023年33-34\%的涨幅，可带来3-4\%的收入增量，且大部分直接转化为利润。二季度GMV增速可能超市场预期的27\%，基于shop.app子域名流量与GMV 94\%的相关性。
{\small\color{refgray}数据：非Plus套餐上次涨价33-34\%（2023年）\par}
{\small\color{refgray}数据：Plus计划涨价25\%（2024年）\par}
{\small\color{refgray}数据：新模式下前四年合作伙伴佣金高104-137\%\par}
{\small\color{refgray}数据：十年累计佣金反而低18\%\par}
{\small\color{refgray}边际变化：从商户数量增长转向ARPU提升，合作伙伴激励改革和定价调整是核心驱动力。\par}
\paragraph{Bernstein} 达美航空二季度调整后每股收益1.56美元，落在自身上限1.50美元之上，同时重申全年 美元的指引。这份业绩本身不算意外——市场此前已有预期——但真正值得关注的不是数字本身，而是支撑这些数字的定价环境正在发生什么变化。
{\small\color{refgray}数据：核心依据不是燃油成本下降，也不是需求突然爆发，而是他们观察到的一个更持久的转变：美国航空业的定价行为正在从“周期性的供需波动”转向“结构性的定价纪律”。报告用了一个关键词组——“positive structural pricing shift”。如果这个判断成立，那么达美航空目前的估值可能还没有完全反映这一变化。\par}
{\small\color{refgray}数据：二季度超预期背后，真正起作用的是“非票收入”和“企业客户”\par}
\paragraph{DB} 这份DB研究的核心判断，来自一场与前动力总成采购主管 Antonio Trotta 的专家对话。结论直接且反直觉：在超豪华汽车领域，竞争壁垒不是技术参数，而是品牌地位与客户心智。报告认为，从今天看，没有明显的竞争不确定性能撼动法拉利的独特位置。
{\small\color{refgray}数据：对于关注高端制造与消费升级的读者，这份报告提供了一个观察框架：当一家公司拥有“价格是最后考虑因素”的客户群时，其定价策略、商业模式与竞争格局，与大众市场完全不同。\par}
{\small\color{refgray}数据：品牌护城河来自 F1 与公路车的双重遗产，竞争对手难以复制\par}
\subsubsection*{关键数据}
\begin{itemize}
  \item 西欧美妆个护价值增速+4.1\%，销量增速+4.7\%
  \item 京东即时零售亏损从135亿收窄至60亿
  \item 中国宠物食品抖音GMV+34\%，天猫淘宝+3\%
  \item 欧洲奢侈品有机增长预计4.7\%，高于市场预期4.3\%
  \item 洽洽Q2净利润超预期12\%，隐含利润率6.4\%
  \item 日本便利店零售媒体市场预计2029年达1.3万亿日元
  \item Shopify非Plus套餐若涨价33-34\%可带来3-4\%收入增量
\end{itemize}
\begin{figure}[htbp]
\centering
\includegraphics[width=0.88\linewidth]{figures/fig_003.jpg}
\caption{FIGURE 1 - \#\# Setting the scene: what are we expecting for Q2? Looking at 2Q26, we expect our coverage to grow \$4.7\textbackslash{}\%\$ on an organic basis, slightly above BBG cons. at \$4.3\textbackslash{}\%\$ and accelerating vs 1Q26 at \$3.5\textbackslash{}\%\$ . Our estimates range from Richemont Jewellery Maisons at \$}
\end{figure}
\begin{figure}[htbp]
\centering
\includegraphics[width=0.88\linewidth]{figures/fig_035.jpg}
\caption{EXHIBIT 2 - EXHIBIT 2: Food mkt value/volume/price growth Notes: Growth rates are weighted by our company sales in the categories EXHIBIT 3: Food Private Label (PL) market share change Note: PL market share is weighted by our company sales}
\end{figure}
\begin{figure}[htbp]
\centering
\includegraphics[width=0.88\linewidth]{figures/fig_100.jpg}
\caption{Figure 1 - Figure 1: Sentiment has not recovered back to pre-Middle East levels European luxury sector YTD 12m Fwd P/E change Today Change Pre-Middle East tension Beginning of year}
\end{figure}
\begin{figure}[htbp]
\centering
\includegraphics[width=0.88\linewidth]{figures/fig_126.jpg}
\caption{Exhibit 1 - Exhibit 1: Category growth on Tmall, Taobao}
\end{figure}
{\small\color{refgray}本节综合 10 篇研究；涉及 BARC、Bernstein、NOM、DB、GS、JEF。}
\newpage
\subsection{医疗与教育：分化加速，高端化与创新驱动结构性机会}
\textbf{中国医药行业在反腐与集采常态化下，创新药企分化加速，GLP-1赛道线上销售强劲；全球CRO行业基本面改善，订单与资金流回暖，AI替代担忧被修正；中国教育市场K-12学生总数下降，但学习平板均价上涨，高端化趋势明显。}
\subsubsection*{主流共识}
\begin{itemize}
  \item 中国医药行业反腐与集采对头部创新药企的实质影响正在收窄，集采从行业级冲击退化为产品级调整。
  \item 全球CRO行业下游客户资金面显著改善，2026年第二季度生物技术股权融资创历史第三高，订单与并购活动回暖。
  \item 中国K-12学生总数2025年同比下降2\%，预计2025-2030年以年均4\%递减，但初高中学生数仍维持正增长。
\end{itemize}
\subsubsection*{分歧与条件差异}
\begin{itemize}
  \item GLP-1线上销售增长强劲，但价格竞争在年初调整后二季度趋于稳定，市场对后续价格走势存在分歧。
  \item CRO板块买方预期高于卖方共识，但持仓仍偏低，实际订单数据若落在买方预期区间，当前估值可能未充分定价。
  \item 学习平板销量同比下降15\%，但均价同比上涨26\%，市场从铺量转向提价，高端化能否持续存疑。
\end{itemize}
\subsubsection*{机构观点}
\paragraph{Bernstein} 中国医药行业分化加速，反腐与集采对头部创新药企影响收窄。恒瑞七氟烷集采影响5-13亿元，但创新药占比提升使公司层面风险可控；豪森创新药收入占比近80\%。GLP-1线上销售强劲，玛仕度肽全年线上销售额有望达13-15亿元，支撑信达生物20亿元销售目标。下一代肿瘤治疗手段增长显著，PD-(L)1双抗同比增长238\%，TROP2 ADC同比增长149\%。
{\small\color{refgray}数据：恒瑞七氟烷占营收约7\%，集采影响5-13亿元\par}
{\small\color{refgray}数据：豪森创新药收入占比接近80\%\par}
{\small\color{refgray}数据：2026年上半年GLP-1线上销售额合计接近30亿元\par}
{\small\color{refgray}数据：玛仕度肽全年线上销售额有望达13-15亿元\par}
{\small\color{refgray}边际变化：从行业级冲击转向产品级调整，创新药占比提升使头部企业抗风险能力增强；GLP-1线上销售数据成为消费医疗需求直接窗口。\par}
\paragraph{DB} 全球CRO行业复苏从是否发生进入能持续多久阶段。2026年第二季度生物技术股权融资创历史第三高，药物批准数量有望超2025年全年，许可交易金额同比增长70\%。买方对订单预期高于卖方共识0.05倍，但持仓仍低配。AI替代担忧被修正，Top 5药企计划长期维持外包比例，AI用于提升质量而非削减成本。定价环境趋稳，FSP和FSO预计持平或小幅调整。
{\small\color{refgray}数据：CRO公司近一个月平均上涨17\%\par}
{\small\color{refgray}数据：2026年第二季度生物技术股权融资历史第三高\par}
{\small\color{refgray}数据：许可交易金额同比增长70\%\par}
{\small\color{refgray}数据：买方book-to-bill预期高于卖方共识0.05倍\par}
{\small\color{refgray}边际变化：从超跌反弹转向基本面改善驱动，AI替代担忧导致的估值折价仍有1-2倍修复空间。\par}
\paragraph{GS} 中国K-12学生总数2025年同比下降2\%，小学和学前教育分别下滑4\%和10\%，初高中维持2\%和4\%正增长。非学科培训机构数量6月环比再降0.2\%。学习平板市场销量承压，六大品牌合计GMV同比下降15\%，但均价同比上涨26\%至约4100元，好未来T6系列定价6900-12000元，高端化趋势明显。学而思智能学习设备App月活跃用户同比增长73\%至约390万。
{\small\color{refgray}数据：2025年K-12学生总数同比下降2\%\par}
{\small\color{refgray}数据：小学和学前教育分别下滑4\%和10\%\par}
{\small\color{refgray}数据：初高中学生数维持2\%和4\%正增长\par}
{\small\color{refgray}数据：2025-2030年K-12学生数年均复合变化率-4\%\par}
{\small\color{refgray}边际变化：从铺量转向提价，高端化趋势测试家长对内容和服务溢价的接受度。\par}
\subsubsection*{关键数据}
\begin{itemize}
  \item 恒瑞七氟烷集采影响5-13亿元
  \item 2026年上半年GLP-1线上销售额合计接近30亿元
  \item PD-(L)1双抗同比增长238\%
  \item CRO公司近一个月平均上涨17\%
  \item 2026年第二季度生物技术股权融资历史第三高
  \item 2025年K-12学生总数同比下降2\%
  \item 好未来学习平板均价同比上涨26\%至约4100元
  \item 学而思智能学习设备App月活跃用户同比增长73\%至约390万
\end{itemize}
\begin{figure}[htbp]
\centering
\includegraphics[width=0.88\linewidth]{figures/fig_008.jpg}
\caption{EXHIBIT 2 - EXHIBIT 3: Mazdutide recorded over CNY 600 Mn online sales in 1H26, getting 22\% of share as 2nd place among Top 3 GLP1s 26H1 GLP-1 online sales by product Semaglutide: include ozempic and wegovy}
\end{figure}
\begin{figure}[htbp]
\centering
\includegraphics[width=0.88\linewidth]{figures/fig_108.jpg}
\caption{Exhibit 1 - Exhibit 1: Number of K-12 students in China peaked out in 2022 and declined -2\% yoy in 2025 Exhibit 2: Primary school and pre-school student count declined -4\% yoy/-10\% yoy, while junior secondary and high school student count ma}
\end{figure}
\begin{figure}[htbp]
\centering
\includegraphics[width=0.88\linewidth]{figures/fig_009.jpg}
\caption{EXHIBIT 3 - EXHIBIT 3: Mazdutide recorded over CNY 600 Mn online sales in 1H26, getting 22\% of share as 2nd place among Top 3 GLP1s 26H1 GLP-1 online sales by product Semaglutide: include ozempic and wegovy \#\# 1Q26 sample hospital sales Kelu}
\end{figure}
\begin{figure}[htbp]
\centering
\includegraphics[width=0.88\linewidth]{figures/fig_109.jpg}
\caption{Exhibit 1 - Exhibit 3: Breakdown of after-school tutoring licenses (Jun 2026)}
\end{figure}
{\small\color{refgray}本节综合 3 篇研究；涉及 Bernstein、DB、GS。}
\newpage
\subsection{工业与基础设施：周期分化与结构性机遇}
\textbf{全球工业与基础设施板块呈现显著分化：AI算力需求驱动电力设备与备电系统结构性增长，铝价上行验证资源端定价权，而传统工程机械面临库存回补与终端需求脱节的低能见度阶段。日本交运板块短期弹性与中期隐忧并存，中国电网设备订单强劲但收入确认节奏扰动短期业绩。}
\subsubsection*{主流共识}
\begin{itemize}
  \item AI机柜功率密度持续攀升，机柜级备电电池（BBU）正从可选项变为必选项，3kW/5kW产品持续放量，8kW/12kW模块预计2026年Q4开始贡献收入。
  \item 铝价上涨是铝业公司利润超预期的主要驱动，中国宏桥1H26税后利润同比增39\%至188亿元，略超预期，强劲现金流有望支撑分红与回购。
  \item 中国电网设备需求基本面健康，思源电气新订单增速达35\%，对美变压器出货量从每月2-3台增至7月约10台，但收入确认延迟和外汇损失造成短期业绩低于预期。
\end{itemize}
\subsubsection*{分歧与条件差异}
\begin{itemize}
  \item BBU市场：Citi认为PSU供应商难以垂直整合进入BBU，台系供应商有望从松下手中夺取份额；而市场担忧HVDC商业化推迟可能冲击BBU主题，Citi认为近端需求主要来自现有架构功率提升。
  \item 北美工程机械：GS判断当前出口强劲（液压挖掘机对北美出口同比增超40\%）主要来自经销商补库存，而非终端需求全面复苏；补库完成后若终端未跟上，当前强势可能为阶段性峰值。
  \item 日本航运：短期看中东局势和财报上调，但2027年后集装箱船交付增加可能松动供需；汽车运输船租金和自有船资产价值被低估，可能成为中期估值催化剂。
  \item 中国AI大模型融资：JPM认为同一笔融资对不同公司含义相反——智谱AI融资加速已可见的收入转化路径，而MiniMax面临确定的稀释成本和尚未证明的转化路径。
\end{itemize}
\subsubsection*{机构观点}
\paragraph{Citi} BBU竞争壁垒正在变高，PSU供应商难以垂直整合进入该市场，台系供应商如Dynapack有望从松下手中夺取份额。随着功率从3kW/5kW向8kW/12kW乃至25kW演进，电池包设计、BMS集成、热管理等能力要求同步提升。HVDC商业化推迟不会破坏BBU逻辑，近端需求主要来自现有架构功率提升，8kW/12kW BBU预计2026年Q4开始贡献收入，25kW模块验证在2027年上半年。
{\small\color{refgray}数据：BBU功率从3kW/5kW向8kW/12kW乃至25kW演进\par}
{\small\color{refgray}数据：8kW/12kW BBU预计2026年Q4开始贡献收入\par}
{\small\color{refgray}数据：25kW模块产品验证预计2027年上半年进行\par}
{\small\color{refgray}数据：Dynapack与PSU厂商的绑定关系构成结构性优势\par}
{\small\color{refgray}边际变化：市场此前认为HVDC商业化是BBU采用必要条件，Citi指出近端需求主要来自现有架构功率提升，HVDC是长期增量而非前提。\par}
\paragraph{Citi} 中国宏桥1H26初步税后利润同比增39\%至188亿元，略超市场预期，主要归因于铝价上涨。1H26净利润已覆盖Bloomberg一致预期的49\%和Citi全年预测的52\%，环比增长61\%显示铝价上涨节奏在加速。强劲盈利和现金流有望在市场波动中支撑分红与股份回购。
{\small\color{refgray}数据：1H26税后利润同比增39\%至188亿元\par}
{\small\color{refgray}数据：1H26净利润约166亿元，环比增长61\%\par}
{\small\color{refgray}数据：1H26净利润覆盖全年预测的52\%\par}
{\small\color{refgray}数据：目标价48港元，基于13倍2026年预期市盈率\par}
{\small\color{refgray}边际变化：铝价上涨的传导效率比市场此前预期更直接、更完整，成本端控制与产能利用率均处于有利区间。\par}
\paragraph{GS} 思源电气二季度营收6235百万元，同比增18\%，但低于预期约8\%，主因海外合同收入确认延迟和外汇损失超1亿元，而非需求转弱。新订单增速仍达35\%，对美变压器出货量从每月2-3台增至7月约10台，预计三季度确认。毛利率已环比回升，下半年北美变压器和特高压项目将开始确认收入。下调2026-2030年EPS 8-11\%，其中3-6\%来自收入延迟，1-2\%来自经营杠杆下降，约3\%来自外汇损失。
{\small\color{refgray}数据：二季度营收6235百万元，同比增18\%，低于预期约8\%\par}
{\small\color{refgray}数据：新订单增速35\%，对美变压器出货量从2-3台/月增至10台/月\par}
{\small\color{refgray}数据：上半年外汇损失超1亿元，去年同期为收益\par}
{\small\color{refgray}数据：下调2026-2030年EPS 8-11\%\par}
{\small\color{refgray}边际变化：市场可能将业绩低于预期解读为需求走弱，GS明确指出是收入确认节奏和汇率问题，订单和出货量保持健康。\par}
\paragraph{GS} 中国机构资管二季度盈利强劲，九家已披露业绩预告的机构平均净利润同比增约155\%，CICC预计增78\%-90\%，中信增70\%。A股二季度日均成交额创纪录达2.9万亿元，两融余额突破3万亿元同比增63\%，IPO融资额660亿元同比增206\%。市场真正关注的是三季度能否将交易量和IPO热度转化为ROE持续修复，IPO项目储备而非情绪是成交额持续性的关键。
{\small\color{refgray}数据：九家机构1H26平均净利润同比增约155\%\par}
{\small\color{refgray}数据：A股2Q26日均成交额创纪录达2.9万亿元\par}
{\small\color{refgray}数据：两融余额突破3万亿元，同比增63\%\par}
{\small\color{refgray}数据：A股IPO融资额660亿元，同比增206\%\par}
{\small\color{refgray}边际变化：市场此前关注交易量持续性，GS指出IPO项目储备才是更可追踪的前瞻指标，而非散户入市节奏。\par}
\paragraph{GS} 日本交运板块报价已部分反映正常化预期，但结构性变化尚未被广泛定价。航运短期弹性来自中东局势和财报上调，日本邮船FY3/27每股股息预期210日元高于指引200日元，但2027年后集装箱船交付增加可能松动供需。汽车运输船租金和自有船资产价值被低估，可能成为中期估值催化剂。造船股回调是资金轮动而非基本面恶化，名村造船和三井E\&S FY3/27 PE约10倍。
{\small\color{refgray}数据：日本邮船FY3/27每股股息预期210日元，高于公司指引200日元\par}
{\small\color{refgray}数据：名村造船和三井E\&S FY3/27 PE约10倍\par}
{\small\color{refgray}数据：集装箱船交付量增加，供需可能从2027年起松动\par}
{\small\color{refgray}数据：汽车运输船租金和自有船资产价值被低估\par}
{\small\color{refgray}边际变化：市场聚焦集装箱船和短期交易逻辑，GS指出汽车运输船和自有船资产重估是尚未被定价的中期催化剂。\par}
\paragraph{GS} 全球工程机械处于结构性低能见度阶段，不同终端市场周期相位各异。北美建筑机械出口强劲（液压挖掘机对北美出口同比增超40\%）主要来自经销商补库存而非终端需求复苏，美国建筑研究过去1-2年基本持平。小型机械中微型挖掘机库存过剩，装载机难以复制过去增长。久保田被列为首选，因其即便终端需求不反弹，利润仍可依靠汇率、关税和成本改善持续增长。
{\small\color{refgray}数据：液压挖掘机对北美出口同比增超40\%\par}
{\small\color{refgray}数据：美国建筑研究过去1-2年基本持平\par}
{\small\color{refgray}数据：微型挖掘机库存过高，需求集中在住房领域\par}
{\small\color{refgray}数据：北美大小型拖拉机零售销量已降至历史低位\par}
{\small\color{refgray}边际变化：市场可能将出口强劲误判为需求复苏，GS明确指出驱动力是经销商补库存，终端需求尚未全面回暖。\par}
\paragraph{JPM} 中国独立大模型厂商2026年至今累计融资超200亿美元，但同一笔融资对不同公司含义相反。智谱AI近期完成约40亿美元配售仅稀释4.4\%，瓶颈已从需求不足转向服务容量不足，新增推理资源可在12个月内直接转化为ARR增长，隐含2027年末ARR 40-50亿美元。MiniMax完成约21亿美元融资稀释约11-17\%，但模型能力尚未证明领先，融资只是争取时间。
{\small\color{refgray}数据：2026年至今中国独立大模型厂商累计融资超200亿美元\par}
{\small\color{refgray}数据：智谱AI配售约40亿美元，仅稀释4.4\%\par}
{\small\color{refgray}数据：智谱AI隐含2027年末ARR 40-50亿美元\par}
{\small\color{refgray}数据：MiniMax融资约21亿美元，稀释约11-17\%\par}
{\small\color{refgray}边际变化：市场普遍认为融资利好所有公司，JPM指出融资是放大镜——放大已有业务能力，而非凭空创造竞争力。\par}
\subsubsection*{关键数据}
\begin{itemize}
  \item AI机柜BBU功率从3kW/5kW向8kW/12kW乃至25kW演进
  \item 中国宏桥1H26税后利润同比增39\%至188亿元
  \item 思源电气新订单增速35\%，对美变压器出货量从2-3台/月增至10台/月
  \item A股2Q26日均成交额创纪录达2.9万亿元
  \item A股IPO融资额660亿元，同比增206\%
  \item 液压挖掘机对北美出口同比增超40\%
  \item 中国独立大模型厂商2026年至今累计融资超200亿美元
  \item 智谱AI配售40亿美元仅稀释4.4\%，隐含2027年末ARR 40-50亿美元
\end{itemize}
\begin{figure}[htbp]
\centering
\includegraphics[width=0.88\linewidth]{figures/fig_113.jpg}
\caption{Exhibit 1 - Exhibit 1: Sieyuan's shipments to the US accelerated to 10 units in July, which will more likely be recognized in 3Q26E Exhibit 2: ... where we estimate transformers (likely c.300MVA) to be priced at c.US\textbackslash{}\$10mn/c.Rmb70mn each E}
\end{figure}
\begin{figure}[htbp]
\centering
\includegraphics[width=0.88\linewidth]{figures/fig_131.jpg}
\caption{Exhibit 1 - Exhibit 3: Hydraulic excavator export volume to North America is driving growth at over 40\% yoy YTD Shipment statistics: Hydraulic excavator shipment volume to North America (yoy)}
\end{figure}
\begin{figure}[htbp]
\centering
\includegraphics[width=0.88\linewidth]{figures/fig_154.jpg}
\caption{Figure 1 - Figure 1: Zhipu \& MiniMax 2022-25 R\&D spending and IPO proceeds 6-month spending (US\textbackslash{}\$mn) Figure 2: Zhipu \& MiniMax 2022-25 R\&D spending as \% of revenue \#\# Zhipu AI: financing improves execution visibility across model develop}
\end{figure}
\begin{figure}[htbp]
\centering
\includegraphics[width=0.88\linewidth]{figures/fig_117.jpg}
\caption{Exhibit 2 - Exhibit 4: The balance of margin financing and securities lending surpassed Rmb 3tn, a YoY increase of 63\%, reflecting a continuous improvement in investor risk appetite. Exhibit 5: A-share IPO fundraising volume reached Rmb 66bn}
\end{figure}
{\small\color{refgray}本节综合 7 篇研究；涉及 Citi、GS、JPM。}
\newpage
\subsection{能源与大宗商品：霍尔木兹扰动长期效力递减，中国需求结构性支撑}
\textbf{短期地缘冲突推升油价脉冲，但中东管道扩建正结构性削弱霍尔木兹海峡的长期封锁风险；同时，中国在半导体测试设备等高端制造领域的需求强劲，为相关大宗商品及设备市场提供独立支撑。}
\subsubsection*{主流共识}
\begin{itemize}
  \item 中东绕过霍尔木兹海峡的管道运力将在2028年底前显著增加，有效降低对海峡的依赖。
  \item AI/HPC芯片测试设备规格持续升级，推动平均售价（ASP）两位数增长。
  \item 中国市场需求保持强劲，为全球供应链提供稳定增量。
\end{itemize}
\subsubsection*{分歧与条件差异}
\begin{itemize}
  \item 管道建设速度存在分歧：GS认为单国项目中位建设时间仅2.5年，但多国协调项目可能拖慢进度，保守与加速情景下覆盖范围相差30个百分点。
  \item 短期地缘溢价持续性存疑：部分观点认为油轮遇袭和军事打击将维持油价高位，而GS强调长期管道扩张会稀释每次冲击的效力。
  \item 中国需求的结构性驱动来源不同：JPM聚焦于AI芯片测试设备，而其他机构可能更关注传统能源或基础金属。
\end{itemize}
\subsubsection*{机构观点}
\paragraph{GS} 中东管道扩建正结构性削弱霍尔木兹海峡的封锁风险。到2028年底，有效管道运力将超过1400万桶/日，覆盖波斯湾七国战前出口量的60\%以上。历史数据显示，中东单国管道中位建设时间仅2.5年，供应中断会加速建设。基准情景下，2027年底新增有效运力380万桶/日，2028年底累计增加730万桶/日。加速情景可覆盖75\%出口，保守情景也能覆盖45\%。近期冲突已将长期油价预期上调9美元至76美元，但管道扩张将限制未来地缘溢价的幅度。
{\small\color{refgray}数据：到2028年底有效管道运力超1400万桶/日\par}
{\small\color{refgray}数据：覆盖波斯湾七国战前出口量60\%以上\par}
{\small\color{refgray}数据：2027年底新增有效运力380万桶/日\par}
{\small\color{refgray}数据：2028年底累计增加730万桶/日\par}
{\small\color{refgray}边际变化：从关注短期地缘脉冲转向长期基础设施扩张，强调管道运力增长将系统性降低霍尔木兹封锁的长期效力。\par}
\paragraph{JPM} 中国需求强劲，尤其在AI芯片测试设备领域。鸿海精密2Q26营收超预期12-13\%，达新台币138亿元，同比增长101\%。设备收入同比增长约一倍，但出货量增长远低于此，反映ASP两位数提升。下半年需求来自ASIC、特斯拉AI芯片、AMD Venice CPU及中国市场四轮驱动。设备产能仍紧张，公司持续评估新生产基地。SEMICON Taiwan将展示三款新设备平台，包括大型封装handler（3Q26出货，ASP两位数提升）和CPO光电共测handler（4Q26试产）。
{\small\color{refgray}数据：2Q26营收新台币138亿元，同比增长101\%\par}
{\small\color{refgray}数据：超预期12-13\%\par}
{\small\color{refgray}数据：设备收入同比增长约一倍\par}
{\small\color{refgray}数据：ASP两位数提升\par}
{\small\color{refgray}边际变化：从关注产能利用率转向ASP驱动增长，强调设备规格升级（超高TDP、多区域温控、AOI集成）是利润核心。\par}
\subsubsection*{关键数据}
\begin{itemize}
  \item 到2028年底中东管道有效运力超1400万桶/日，覆盖出口60\%以上
  \item 单国管道中位建设时间2.5年
  \item 2027年底新增有效运力380万桶/日，2028年底累计730万桶/日
  \item 加速情景覆盖75\%，保守情景覆盖45\%
  \item 长期油价预期上调9美元至76美元
  \item 鸿海精密2Q26营收新台币138亿元，同比增101\%，超预期12-13\%
  \item 设备ASP两位数提升，出货量增长远低于收入增长
  \item 大型封装handler 3Q26出货，ASP两位数提升
\end{itemize}
\begin{figure}[htbp]
\centering
\includegraphics[width=0.88\linewidth]{figures/fig_141.jpg}
\caption{Exhibit 3 - Exhibit 1: Our Base Case Implies Over 14mb/d of Effective Pipeline Capacity by End-2028, Insulating More Than 60\% of the Pre-War Level of Persian Gulf Producers' Exports From Any Potential Closures}
\end{figure}
\begin{figure}[htbp]
\centering
\includegraphics[width=0.88\linewidth]{figures/fig_142.jpg}
\caption{Exhibit 1 - Exhibit 2: Across a Sample of 9 Major Pipelines, the Median Existing Pipeline Was Built in 2.5 Years Encompasses eight pipelines built in the Persian Gulf and the fastest pipeline ever built (US + Canada). The East-West Pipeline}
\end{figure}
\begin{figure}[htbp]
\centering
\includegraphics[width=0.88\linewidth]{figures/fig_143.jpg}
\caption{Exhibit 1 - Exhibit 3: Depending on the Scope of Pipeline Projects and Speed of Their Development, We Estimate Anywhere Between 45–75\% of Pre-War Persian Gulf Producers' Exports Will Be Insulated by Effective Pipeline Capacity By End-2028}
\end{figure}
\begin{figure}[htbp]
\centering
\includegraphics[width=0.88\linewidth]{figures/fig_144.jpg}
\caption{Exhibit 2 - Exhibit 3: Depending on the Scope of Pipeline Projects and Speed of Their Development, We Estimate Anywhere Between 45–75\% of Pre-War Persian Gulf Producers' Exports Will Be Insulated by Effective Pipeline Capacity By End-2028}
\end{figure}
{\small\color{refgray}本节综合 2 篇研究；涉及 GS、JPM。}
\newpage
\subsection{支付处理商估值：低PE背后的“清洁盈利”真相}
\textbf{市场对FIS、Fiserv和Global Payments的低PE估值存在系统性误读，非GAAP盈利因持续计入重组与转型费用而高估真实盈利能力，基于“清洁盈利”框架的真实估值（EV/NOPAT）显著高于表面倍数。}
\subsubsection*{主流共识}
\begin{itemize}
  \item 三家支付处理商当前NTM市盈率处于历史低位（FIS 6.4x、Fiserv 6.0x、GPN 5.1x），吸引价值投资者关注。
  \item 非GAAP盈利调整中，重组、转型和并购整合费用被归类为“一次性”，但实际每年持续发生。
  \item 低税率（如FIS调整后税率12\%-13\%）进一步夸大了调整后每股收益，若按正常20\%税率计算，真实盈利更低。
\end{itemize}
\subsubsection*{分歧与条件差异}
\begin{itemize}
  \item 市场对“调整后”盈利的可靠性存在分歧：部分投资者认为低PE反映真实价值，另一部分质疑会计调整的可持续性。
  \item 对于Fiserv，其清洁盈利占非GAAP比例（94\%）显著高于FIS（77\%）和GPN（75\%），表明公司间调整质量差异明显。
  \item 回购能力分化：GPN因并购交易性支出导致清洁FCF转化率低于同行，影响股东回报预期。
\end{itemize}
\subsubsection*{机构观点}
\paragraph{Bernstein} Bernstein通过构建“清洁盈利”框架，重新计算FIS、Fiserv和Global Payments的真实估值。2025年清洁盈利占非GAAP盈利比例分别为77\%、94\%和75\%，2026年Q1进一步降至76\%、75\%和71\%，表明非GAAP EPS高估真实盈利能力20\%-30\%。调整后EV/NOPAT倍数（2026E）分别为14.3x、10.8x和11.3x，远高于表面PE（6.7x、6.3x、5.6x）。差距主要来自持续发生的重组费用（如Fiserv的“One Fiserv”转型费用超1.4亿美元）和整合成本（GPN的Genius项目重组9600万美元）。低税率（FIS调整后税率12\%-13\%）若正常化至20\%，真实收益将进一步折扣。
{\small\color{refgray}数据：FIS 2025年清洁盈利/非GAAP盈利=77\%\par}
{\small\color{refgray}数据：Fiserv 2025年清洁盈利/非GAAP盈利=94\%\par}
{\small\color{refgray}数据：GPN 2025年清洁盈利/非GAAP盈利=75\%\par}
{\small\color{refgray}数据：Fiserv 2026Q1“One Fiserv”转型费用加回>1.4亿美元\par}
{\small\color{refgray}边际变化：此前市场仅关注低PE倍数，Bernstein首次系统量化“清洁盈利”与非GAAP盈利的差距，揭示真实估值被系统性低估数倍，并指出低税率假设的脆弱性。\par}
\subsubsection*{关键数据}
\begin{itemize}
  \item FIS 2025年清洁盈利/非GAAP盈利=77\%
  \item Fiserv 2025年清洁盈利/非GAAP盈利=94\%
  \item GPN 2025年清洁盈利/非GAAP盈利=75\%
  \item Fiserv 2026Q1“One Fiserv”转型费用加回>1.4亿美元
  \item GPN 2026Q1 Genius项目重组加回9600万美元
  \item FIS调整后税率12\%-13\%，正常化税率20\%
  \item FIS 2026E EV/NOPAT=14.3x vs 表面PE 6.7x
  \item Fiserv 2026E EV/NOPAT=10.8x vs 表面PE 6.3x
  \item GPN 2026E EV/NOPAT=11.3x vs 表面PE 5.6x
\end{itemize}
\begin{figure}[htbp]
\centering
\includegraphics[width=0.88\linewidth]{figures/fig_093.jpg}
\caption{EXHIBIT 6 - EXHIBIT 6: ...they are several turns higher on clean EV/NOPAT (FIS 14.3x, FISV 10.8x, GPN 11.3x) EV/NOPAT Multiple (2026E) EXHIBIT 7: We note that 2027 analysis gets a bit trickier because we have to make a fair bit of assumption}
\end{figure}
\begin{figure}[htbp]
\centering
\includegraphics[width=0.88\linewidth]{figures/fig_090.jpg}
\caption{EXHIBIT 2 - EXHIBIT 4: We have been closely watching ‘clean’ FCF as \% of reported adj. FCF Clean FCF Conversion (as \% of Adj. FCF) For GPN: we added back one-time transaction-related costs due to the WP deal closing to clean FCF}
\end{figure}
\begin{figure}[htbp]
\centering
\includegraphics[width=0.88\linewidth]{figures/fig_091.jpg}
\caption{EXHIBIT 4 - EXHIBIT 4: We have been closely watching ‘clean’ FCF as \% of reported adj. FCF Clean FCF Conversion (as \% of Adj. FCF) For GPN: we added back one-time transaction-related costs due to the WP deal closing to clean FCF EXHIBIT 5: W}
\end{figure}
\begin{figure}[htbp]
\centering
\includegraphics[width=0.88\linewidth]{figures/fig_092.jpg}
\caption{EXHIBIT 5 - EXHIBIT 7: We note that 2027 analysis gets a bit trickier because we have to make a fair bit of assumptions on ‘adjustments’. Our estimates for 2027 PE multiples include 7.3x for FIS, 6.3x for FISV, and 5.8x for GPN}
\end{figure}
{\small\color{refgray}本节综合 1 篇研究；涉及 Bernstein。}
\newpage
\subsection{AI数据中心：运营商自建算力服务，收入结构加速分化}
\textbf{中国移动AI数据中心（AIDC）收入增速远超传统数据中心，自建自营模式与AI运维能力正推动运营商从资源出租向算力服务转型，成为AI基础设施浪潮中的结构性受益者。}
\subsubsection*{主流共识}
\begin{itemize}
  \item AI算力需求正从训练向推理延伸，推动数据中心收入结构分化，AIDC增速显著高于传统业务。
  \item 运营商自建自营模式在供电、冷却、运维等方面更具适配性，能深度参与客户AI工作流。
  \item AI技术本身（如智能调度、数字孪生）正被用于提升数据中心运营效率。
\end{itemize}
\subsubsection*{分歧与条件差异}
\begin{itemize}
  \item 部分机构认为运营商AIDC增长主要依赖国内需求，而GS调研强调海外中企出海也将贡献增量。
  \item 关于AIDC利润率能否持续高于传统数据中心，市场存在分歧，GS认为自建模式可维持优势。
  \item 部分观点担忧电力与冷却成本上升侵蚀利润，但GS指出前瞻性设计已考虑长期可持续性。
\end{itemize}
\subsubsection*{机构观点}
\paragraph{GS} 中国移动香港全球智算中心自建自营，提供算力、网络、供电、冷却一体化服务，直接支持大模型训练与推理。2025年AIDC收入同比增长35\%，而整体数据中心收入仅增9\%，AI算力已成为核心增长引擎。中心自身运用AI进行智能调度、机器人巡检与数字孪生，运营效率提升明显。设计已预判算力增长与电力消耗趋势，长期可持续性优于合作模式。海外业务伴随中企出海同步拓展，有望持续推动AI收入占比提升。
{\small\color{refgray}数据：2025年AIDC收入同比+35\%，整体数据中心收入+9\%\par}
{\small\color{refgray}数据：香港全球智算中心2026年3月启用，为移动在香港第二大数据中心\par}
{\small\color{refgray}数据：目标价94港元，基于2027年EV/EBITDA 4.2倍\par}
{\small\color{refgray}边际变化：从传统数据中心资源出租转向自建自营算力服务，AI运维技术从概念进入实际部署，AIDC收入增速与整体数据中心差距拉大至26个百分点。\par}
\subsubsection*{关键数据}
\begin{itemize}
  \item 2025年AIDC收入同比+35\%，整体数据中心收入+9\%
  \item 香港全球智算中心2026年3月启用，为移动在香港第二大数据中心
  \item 目标价94港元，基于2027年EV/EBITDA 4.2倍
\end{itemize}
{\small\color{refgray}本节综合 1 篇研究；涉及 GS。}
\newpage
\section{辅助信号｜战略咨询}
本板块涵盖波士顿咨询两份报告：印度矿业与工程机械的结构性崛起，以及AI在医疗患者接入中心的价值重塑。前者强调印度MCE行业从追赶者到净出口国的转变，后者指出AI应被用于捕捉需求信号而非仅降本。
\subsection{印度矿业与工程机械的结构性增长与出口跃升}
\textbf{印度MCE行业正从依赖进口的小市场转变为全球增长最快的平台之一，未来五年全球份额将从4\%升至约6.5\%，出口目标至2047年或超750亿美元。增长由基础设施扩张、矿业增产、本地化制造深化和出口能力提升四大结构性驱动力支撑，而非短期政策刺激。}
\begin{itemize}
  \item 印度MCE国内需求2025年超170亿美元，全球市场约4200-4400亿美元，亚太占主导。
  \item 基础设施持续扩张：国家高速公路网从2014年9.1万公里扩至2025年约14.6万公里；地铁运营里程从248公里增至超1000公里；运营机场数量翻倍至约160个。
  \item 矿业产量十年跃升：煤炭从6.09亿吨增至10亿吨，铁矿石从1.55亿吨增至3亿吨，石灰石从3.2亿吨增至4.5亿吨，推动下游钢铁、水泥和电力需求。
  \item 印度MCE出口十年翻近3倍至49亿美元，正成为净出口国；BEML案例显示本土化率超90\%，75\%以上业务来自自主研发产品。
  \item 全球MCE行业正经历清洁动力总成、数字化车队、自动化、地下采矿机械化等六重技术变革，印度企业能否同步升级决定其全球竞争力上限。
\end{itemize}
\begin{figure}[htbp]
\centering
\includegraphics[width=0.88\linewidth]{figures/fig_197.jpg}
\caption{EXHIBIT 2 - 23. UN Comtrade data (HS-code trade flows), 2024 24. DGCIS trade data (HS-code trade flows), 2025 \#\# EXHIBIT 2 \#\# India's exports in Mining and Construction Equipment sector have been growing over the years \#\# EXHIBIT 3 \#\# India holds \$\textbackslash{}sim 4\textbackslash{}\%\$ of MCE import }
\end{figure}
\begin{figure}[htbp]
\centering
\includegraphics[width=0.88\linewidth]{figures/fig_198.jpg}
\caption{EXHIBIT 3 - \#\# EXHIBIT 3 \#\# India holds \$\textbackslash{}sim 4\textbackslash{}\%\$ of MCE import demand across major regions Import Value (USD billions) \# Two emerging opportunity clusters for Indian players - each with its own requirement \#\# Cluster 1: Southeast Asia, Africa, the Gulf and South America}
\end{figure}
\subsection{AI在患者接入中心：从成本削减到利润引擎}
\textbf{医院若仅将AI视为降本工具，将错失3-5倍的利润提升机会。真正的价值在于将每次来电转化为需求信号、患者关系和照护缺口关闭。超过50\%的来电可由AI完全自主处理，但成功关键在于内部业务逻辑标准化。}
\begin{itemize}
  \item 目前高达40\%的预约相关交互在首次接触时未能解决，大量需求因等待、转接失败或非工作时间流失。
  \item AI可自主处理高频、脚本化、非临床交互（预约确认、日程安排、取消、转诊查询、导诊问询），占大型医疗系统总来电量的50\%以上。
  \item 无单一厂商提供完整端到端虚拟座席方案，医院需在购买成熟产品与自建系统间权衡，评估技术基础设施、业务逻辑复杂度、供应商依赖和内部工程能力。
  \item 未来token成本变动是不确定变量，需纳入财务模型；内部预约规则、保险验证逻辑等基础数据若不标准化，AI无法可靠运行。
  \item AI赋能后，人工座席可转向高价值患者互动，将成本中心转化为收入引擎，实现3-5倍利润提升。
\end{itemize}
\begin{figure}[htbp]
\centering
\includegraphics[width=0.88\linewidth]{figures/fig_200.jpg}
\caption{Exhibit 1 - Many of our clients are deploying AI to automate lower-complexity processes while reorienting hardworking human staff to high-value, high-impact consumer engagement. They understand that the access center of the future should be an AI-enabled platform to drive}
\end{figure}
\begin{figure}[htbp]
\centering
\includegraphics[width=0.88\linewidth]{figures/fig_201.jpg}
\caption{Exhibit 2 - \# Automate and fully contain end-to-end access workflows. Across health systems, the interactions that consistently demonstrate the highest potential for end-to-end automation are appointment information and confirmation, scheduling and rescheduling, and cance}
\end{figure}
\newpage
\section{辅助信号｜智库/国际机构}
本板块聚焦国际清算银行最新研究，揭示AI对私人信贷软件敞口和全球资本开支竞赛的双重冲击。核心信号：软件贷款风险未被定价，AI投资规模已超社会最优水平1.5-3倍，融资结构脆弱性可能通过循环持股放大。
\subsection{AI颠覆私人信贷中的软件企业敞口}
\textbf{生成式AI正从产品替代和成本压缩两个方向冲击软件企业现金流，但私人信贷市场尚未对此进行差异化定价，形成潜在的系统性风险。}
\begin{itemize}
  \item 美国商业发展公司对软件企业贷款总额约1150亿美元，占科技组合80\%以上，但信用利差与其他行业趋同，未反映AI带来的不确定性。
  \item 软件曾是现金流最可靠的行业（订阅制、高毛利、低资本需求），但生成式AI可替代产品、降低开发门槛、加剧竞争，直接削弱借款企业现金流确定性。
  \item 截至2025年底，软件贷款逾期率不足1\%，但滞后指标可能掩盖风险；非上市BDC贷款利差更低，市场监督最薄弱。
  \item BDC估值未因软件敞口大小而分化，表明市场尚未对AI冲击进行差异化定价，定价信号可能制造虚假安全感。
\end{itemize}
\subsection{AI投资竞赛：7000亿美元资本开支的脆弱性}
\textbf{全球AI基础设施资本开支预计2026年突破7000亿美元，但赢家通吃的竞赛逻辑导致过度投入约1.5倍社会最优水平，融资结构（杠杆+循环持股）放大下行风险。}
\begin{itemize}
  \item BIS模型显示，当前AI投资规模已超社会最优水平约1.5倍，需求弹性低时可达3倍；竞赛外部性驱动企业投入远超理性水平。
  \item 早期投入不可逆，若AI生产力低于预期，产能过剩和资产减值将非线性放大，且越大的泡沫越容易自我击穿。
  \item 融资脆弱性来自杠杆扩张和循环持股（投资AI实验室换取算力承诺），一家企业承压可能通过股权链条传导至整个体系。
  \item 债务融资的专用资产（如GPU）在危机时难以快速变现，加剧流动性风险。
  \item 模型揭示不对称性：生产力超预期收益线性，不及预期损失非线性且通过资金网络放大。
\end{itemize}
\begin{figure}[htbp]
\centering
\includegraphics[width=0.88\linewidth]{figures/fig_190.jpg}
\caption{Figure 1 - Figure 1: AI boom and its financing (a) AI investment boom and historical episodes (b) AI's circular financing structure Note: Panel (a) shows investment during five innovation-led booms as multiples of each episode's pre-boom}
\end{figure}
\begin{figure}[htbp]
\centering
\includegraphics[width=0.88\linewidth]{figures/fig_192.jpg}
\caption{Figure 2 - Figure 2: Investment, the bust threshold, and output (a) Output along the investment path. (b) The endogenous bust frontier. Note: Panel (a) plots sector output \$\textbackslash{}Theta\$ against period-1 investment \$I\_\{1\}\$ for a favourable dra}
\end{figure}
\section{结语}
后续需验证的关键节点包括：台积电3Q26毛利率能否达到市场预期的69\%-70\%，还是如MS预测被海外厂和2nm稀释至67\%-68\%；日元空头回补是调整尾声还是持续性趋势；中国新能源市场燃油车渗透率反弹是供给端主动去库存的结果还是需求趋势逆转；以及CRO板块订单数据能否落在买方预期区间，验证当前估值是否充分定价。
\newpage
\section{覆盖概览}
投行/券商 68 篇；战略咨询 2 篇；智库/国际机构 2 篇。\par
投行覆盖：MS 17篇 / GS 16篇 / Bernstein 10篇 / JPM 7篇 / NOM 7篇 / DB 4篇 / Citi 3篇 / JEF 2篇 / BARC 1篇 / BofA 1篇\par
\section{Disclaimer}
{\scriptsize\color{refgray}
This community is a paid learning and information-sharing community created for general educational purposes only. The membership fee is charged solely for access to community discussions, educational content, organization, and operational costs. It is not an advisory fee, management fee, performance fee, brokerage fee, or compensation for personalized financial, investment, legal, tax, or professional advice.\par
I participate and share information solely as an individual and for educational discussion among community members. I am not acting as a financial adviser, investment adviser, broker, dealer, portfolio manager, legal adviser, tax adviser, accountant, fiduciary, or any other licensed professional adviser. Nothing shared in this community constitutes, and should not be interpreted as, financial advice, investment advice, legal advice, tax advice, a recommendation, solicitation, offer, endorsement, instruction, or guarantee to buy, sell, hold, short, trade, or invest in any security, cryptocurrency, fund, derivative, strategy, or other financial product.\par
All content shared in this community, including messages, discussions, links, files, charts, examples, market commentary, watchlists, AI-generated or AI-polished summaries, and third-party materials, is general, impersonal, and educational in nature. It does not take into account any individual's financial situation, investment objectives, risk tolerance, investment horizon, liquidity needs, tax status, legal restrictions, or suitability requirements. No content should be relied upon as suitable for any specific person, account, portfolio, or financial decision.\par
Information shared in this community may be incomplete, inaccurate, outdated, speculative, AI-generated, AI-edited, or based on third-party internet sources that have not been independently verified. I make no representation or warranty as to the accuracy, completeness, timeliness, reliability, or suitability of any information shared.\par
Financial markets involve substantial risk, including the possible loss of principal. Past performance, historical data, hypothetical examples, simulated results, backtests, personal opinions, or educational examples do not guarantee future results. Each member is solely responsible for conducting their own research, exercising independent judgment, and making their own decisions. Before making any financial, investment, legal, or tax decision, members should consult qualified and licensed professionals.\par
Participation in this community, payment of any membership fee, or communication with me or other members does not create any adviser-client, fiduciary, professional, contractual, agency, partnership, employment, or other advisory relationship. By joining or participating in this community, you acknowledge and agree that you are solely responsible for your own decisions, actions, transactions, profits, losses, risks, and consequences, and that I shall not be liable for any direct, indirect, incidental, consequential, financial, legal, tax, or other loss or damage arising from or related to any content, discussion, information, or material shared in this community.\par
}
\newpage
\begin{center}
{\Large 更多详情报告kcdesk.com}\par\vspace{0.8cm}
\includegraphics[width=0.58\linewidth]{\detokenize{../../prompts/zsxq_img.jpg}}
\end{center}
\end{document}
